% GCT_KRON_20261120, saturation/obstructions worker.
% All original arguments in this file are findings of this run, not claims of
% historical priority. Known results are identified as such.

\section{Saturation: precise negative and positive statements}

Throughout this part, coefficients and decomposition multiplicities are taken
over $\mathbb C$.  Write $R_{h,d}=(d^h)$ and
$g_{h,d}(\lambda)=g(\lambda,R_{h,d},R_{h,d})$.  Put
$\rho(N)=(N-|\rho|,\rho)$ only when this is a partition.
``Ordinary saturation'' means the implication
\[
 g(t\lambda,t\mu,t\nu)>0\text{ for some }t\geq1
 \quad\Longrightarrow\quad g(\lambda,\mu,\nu)>0.
\]
This must be distinguished from the residue-class saturation index in
Mulmuley's revised GCT VI, Definition~1.2.4 (arXiv:0704.0229v4,
printed p.~14; PDF p.~17), and Hypothesis~3.3.2 (printed p.~51;
PDF p.~54). Identically zero constituents of a quasipolynomial do not
violate that definition.  The revised hypothesis asks for a polynomial
bound on the index, and a generic index-zero statement; it does not assert
ordinary saturation.

\subsection{O-001: exact holes and a factor two for two two-row factors}

The input here is a known theorem, not a new positivity rule: Blasiak,
Mulmuley and Sohoni, GCT IV, Proposition~18.6, equations~(175)--(180),
supplied \texttt{Apr13KroneckerGCT4.tex} version, pp.~128--129 in
the matching author-hosted PDF. The proposition itself has no source
label; the six equation labels are \texttt{e g = 0 1} through
\texttt{e g = 0 6}. The proposition applies to
$\ell(\lambda),\ell(\mu)\leq2$ and $\ell(\nu)\leq4$.
Use column counts
\[
 A=\lambda_1-\lambda_2,\quad B=\mu_1-\mu_2,
 \quad c_i=\nu_i-\nu_{i+1}\ (1\leq i\leq4),\quad\nu_5=0.
\]
All three partitions have the same size.  The known classification says
that the coefficient vanishes if and only if at least one of the following
conditions holds:
\begin{align}
 &(A+B)/2>c_1+c_2+c_3; \tag{O-001.a}\\
 &|A-B|/2>\min(c_1,c_3)+c_2; \tag{O-001.b}\\
 &c_1=c_3=0,\qquad (A+B)/2\not\equiv c_2\pmod2; \tag{O-001.c}\\
 &\min(A,B)=0,\quad \min(c_1,c_2,c_3)=0,
       \quad\max(A,B)/2\not\equiv c_1+c_2\pmod2; \tag{O-001.d}\\
 &\{A,B\}=\{0,2\},\qquad c_1,c_2,c_3\text{ all even};\tag{O-001.e}\\
 &A=B=0,\qquad c_1\text{ or }c_2\text{ odd}. \tag{O-001.f}
\end{align}
The halves in the congruences are integers whenever the other conditions
in that line hold; for instance $A$ and $B$ have the parity of the common
size, and $\min(A,B)=0$ forces both to be even.

\paragraph{Finding O-001.}
In this entire $2\times2\times4$ family, a positive stretch exists if and
only if neither (O-001.a) nor (O-001.b) holds, and this is equivalent to
\[
                         g(2\lambda,2\mu,2\nu)>0.
\]
Consequently, the ordinary-saturation holes are exactly the triples
satisfying the two weak cone inequalities and at least one of
(O-001.c)--(O-001.f).

\begin{proof}
The quantities $A,B,c_i$ all multiply by the stretching factor.  A failure
of either homogeneous cone inequality therefore persists under every
stretch.  At stretch two, (O-001.c), (O-001.d), and (O-001.f) are false
because all relevant integers in their congruences are even.  Condition
(O-001.e) is false: it would require $\{A,B\}=\{0,1\}$ before stretching,
which is impossible because $A$ and $B$ have the same parity.  The known
classification now gives the asserted equivalence and the exact list of
holes.  This proves a factor-two statement, not saturation at factor one.
\end{proof}

For three two-row partitions an even simpler complete description is
available.  Write their second rows as $0\leq c\leq b\leq a\leq N/2$.
Rosas, \emph{The Kronecker product of Schur functions indexed by two-row
shapes or hook shapes}, Corollary~5 (arXiv:math/0001084, p.~11), gives
\[
 g=\#\left\{j\in\mathbb Z:
 \max(0,\lceil(a+b+c-N)/2\rceil)\leq j
                  \leq\lfloor(b+c-a)/2\rfloor\right\}.
\]
Thus a positive stretch exists exactly when $a\leq b+c$; within this cone
the holes are exactly $2a=N$ and $b+c-a$ odd.  To check the last assertion,
put $D=b+c-a\geq0$ and $E=N-2a\geq0$: the interval is
$[\max(0,\lceil(D-E)/2\rceil),\lfloor D/2\rfloor]$.
It is nonempty for $E\geq1$, and for $E=0$ precisely when $D$ is even.
The rectangular parity ray
\[
            g((k,k),(k,k),(k,k))=\begin{cases}1&k\text{ even},\\0&k\text{ odd}
                                      \end{cases}
\]
is already Rosas's Example~6, p.~12. Its revised GCT VI saturation index
is zero.  By contrast, Briand--Orellana--Rosas, Theorem~2.4,
equation~(2.5), gives the strong-saturation counterexamples
\[
\begin{gathered}
 \alpha=(k,k),\quad\beta=(k+1,k-1),\quad
 \gamma=(2k-2i-2j,2i,2j),\quad i>j>0,\ k>2i+j,\\
 g(t\gamma,t\alpha,t\beta)=
 \begin{cases}t/2+1&t\text{ even},\\(t-1)/2&t\text{ odd}.
 \end{cases}
\end{gathered}
\]
Its smallest example is $(\alpha,\beta,\gamma)=((6,6),(7,5),(6,4,2))$.
Here the constituent containing $t=1$ is not identically zero.

\subsection{O-002: exact stretching on long-row rectangular holes}

\paragraph{Finding O-002.}
For $h\geq2$, $d\geq1$, put $\eta=(hd-1,1)$.  For every integer
$t\geq1$,
\[
 g(t\eta,(td)^h,(td)^h)
       =[z^t]\prod_{j=2}^{h}\frac1{1-z^j}.                 \tag{O-002.1}
\]
In particular $g(\eta,R_{h,d},R_{h,d})=0$ and
$g(2\eta,R_{h,2d},R_{h,2d})=1$.  For $h=2$ the support of the ray is
the positive even integers; for $h\geq3$ it is every integer $t\geq2$.
The formal $t=0$ value is one.

\begin{proof}
Use the known rectangular/two-row identity (Pak--Panova, \emph{Strict
unimodality of $q$-binomial coefficients}, Lemma~3, p.~2), in the form
proved in F-002:
\[
 g((htd-t,t),(td)^h,(td)^h)=p_t(h,td)-p_{t-1}(h,td).
\]
The two-row shape is valid because $hd\geq2$. Here $p_s(h,td)$ counts
partitions of $s$ with at most $h$ rows, each of size at most $td$.
Since $td\geq t$, the width bound is inactive for $s=t,t-1$.
Conjugating partitions shows that the generating series is
$\prod_{j=1}^{h}(1-z^j)^{-1}$; multiplication by $1-z$ proves
(O-002.1).  The claimed supports follow from the additive semigroup
generated by $2,\ldots,h$. If $h\geq3$, an even integer $t\geq2$ is a
sum of twos, while an odd integer $t\geq3$ is three plus a sum of twos.
The coefficient at $t=2$ is one.
\end{proof}

This yields holes at arbitrarily large rectangle height, not only at
height two.  If $m\geq2$, the partition $\eta=(hd-1,1)$ satisfies both
coarse necessary padded-permanent conditions
\[
                 \ell(\eta)\leq m^2,
       \qquad \eta_1\geq (h-m)d.
\]
Indeed $\ell(\eta)=2$ and the second inequality is $md\geq1$.
These conditions therefore do not imply ordinary saturation even when
$h$ exceeds any prescribed polynomial in $m$.  They do not make $\eta$
an occurrence candidate: the ambient plethysm coefficient
$a_{\eta}(d[h])$ is zero, as O-005 proves for every nontrivial hook.

\subsection{O-003: a normal, positive, pure geometric model with a hole}

\paragraph{Finding O-003.}
For $h\geq2$, conjugation invariants of a traceless matrix satisfy
\[
 \mathbb C[\mathfrak{sl}_h]^{\mathrm{GL}_h}
       =\mathbb C[c_2,\ldots,c_h],\qquad \deg c_j=j,       \tag{O-003.1}
\]
where $c_j$ are the coefficients of its characteristic polynomial.
This gives a geometric realization of (O-002.1). For $h\geq3$ the
support monoid of the displayed grading is
$\{0,2,3,4,\ldots\}$; its generated group is $\mathbb Z$, and it is
not saturated in that group.  Nevertheless the invariant quotient is a
smooth affine space, its coordinate ring is normal, and multiplication
in the monomial basis has nonnegative structure constants.

\begin{proof}
Let $D$ be the diagonal traceless matrices.  Restriction of invariants to
$D$ is injective: the matrices with pairwise distinct eigenvalues form a
dense open subset of $\mathfrak{sl}_h$, and each is conjugate to a point
of $D$.  Its image lies in $\mathbb C[D]^{S_h}$.
The fundamental theorem of symmetric polynomials, followed by imposing
$e_1=0$, gives
\[
 \mathbb C[D]^{S_h}=\mathbb C[e_2,\ldots,e_h].
\]
For completeness, invariants commute with this quotient: averaging a
polynomial representative over $S_h$ gives a symmetric representative
of each invariant class, and a symmetric polynomial divisible by $e_1$
has symmetric quotient. The functions $e_2,\ldots,e_h$ are algebraically
independent after $e_1=0$ and are restrictions of $c_2,\ldots,c_h$.
This proves (O-003.1) and its graded Hilbert series.

Manivel's rectangular stability theorem (\emph{On rectangular Kronecker
coefficients}, Theorem~1; reproduced as Ikenmeyer--Panova Theorem~2.1)
identifies the coefficient on the left of (O-002.1) with
$\dim(\operatorname{Sym}^t\mathfrak{sl}_h)^{\mathrm{SL}_h}$: transpose
both rectangles if necessary, and use the inequality $td\geq t$.
The trace pairing identifies the adjoint representation with its dual,
and $\mathrm{GL}_h$ and $\mathrm{SL}_h$ have the same conjugation
invariants.  Hence (O-003.1) independently gives the Hilbert series in
(O-002.1).

For $h\geq3$, degrees two and three occur and degree one does not.
Their generated group contains $3-2=1$, so the support is not saturated
even in its own group. A polynomial ring is normal, and its affine
spectrum is smooth. The product of two monomials is a monomial with
coefficient one. These statements prove all the asserted properties.
\end{proof}

\paragraph{Integral and cohomological dictionary.}
The weighted affine scheme
\[
       Q_h=\operatorname{Spec}\mathbb Z[c_2,\ldots,c_h]
\]
and the $\mathbb G_m$-action $c_j\mapsto u^jc_j$ are defined over
$\mathbb Z$. SGA~3, Expos\'e~I, Proposition~4.7.3 (original pp.~35--36)
identifies the grading with this diagonalizable-group action; the
following remark makes the correspondence for algebras and affine
schemes explicit. Its complex generic quotient interpretation is
(O-003.1); no assertion about formation of matrix invariants in all bad
characteristics is needed. For a finite field $\mathbb F_q$ and
$\ell\nmid q$, its underlying affine space has
\[
 H_c^i((Q_h)_{\overline{\mathbb F}_q},\mathbb Q_\ell)=
 \begin{cases}\mathbb Q_\ell(-(h-1))&i=2(h-1),\\0&\text{otherwise}.
 \end{cases}
\]
This follows from SGA~4, Expos\'e~XV, Corollary~2.2, p.~192
(affine-space acyclicity for invertible torsion), and
Expos\'e~XVIII, Theorem~3.2.5, p.~585 (Poincar\'e duality),
followed by the compatible $\ell$-adic limit.
Thus even this pure Tate cohomology does not remove the degree-one hole.
The $\mathbb G_m$-weight of a regular function and a Frobenius weight
on cohomology are different gradings. The same example can be phrased
after base change to a field as $M_c(Q_h)=\mathbb L^{h-1}$ in a
theory of motives with compact supports and affine-space homotopy
invariance.

This is a counterexample to the inference ``normality, positive
structure constants, or purity alone force saturation of degree
support.'' It is not a counterexample to the correct normal-toric
statement about a semigroup algebra: the grading here identifies
distinct monomials, so $\mathbb C[c_2,\ldots,c_h]$ is not the semigroup
algebra of its one-dimensional degree support.

\subsection{O-004: all stretches at least two in a stable rectangular range}

\paragraph{Finding O-004.}
Let $h\geq7$, let $\rho$ be a partition with $\ell(\rho)<h^2$, and
let $d\geq\max(1,|\rho|)$.  Then $\lambda=\rho(hd)$ is a partition and
\[
          g(t\lambda,(td)^h,(td)^h)>0\qquad(t\geq2).       \tag{O-004.1}
\]
For $\rho=\varnothing$ the coefficient is one.  This is a short
strengthening of the even-tail conclusion stated in
Ikenmeyer--Panova Corollary~5.3, obtained from their
Proposition~5.2.  It is not claimed to be a new theorem of the literature.

\begin{proof}
Use $a_\sigma(h)$ for the stable rectangular invariant multiplicity
in Ikenmeyer--Panova Theorem~2.1.  Proposition~5.2 states that for
$h\geq7$, $0\leq k<h^2$, and every $t>1$,
\[
      g\bigl(t(h^2-k,1^k),(th)^h,(th)^h\bigr)>0.
\]
The tail of its first argument is $(t^k)$. By monotonicity as the
rectangle width increases, followed by stability, it follows that
$a_{(t^k)}(h)>0$ for every $k<h^2$.

Write $\rho$ as the rowwise sum of its columns:
$\rho=\sum_{j=1}^{\rho_1}(1^{\rho'_j})$.
The stable multiplicities have the semigroup property. One can see
this without an additional theorem: choose a width large enough that
the two positive coefficients are stable, apply the ordinary
Kronecker semigroup property to their first-row-extended shapes, and
then pass to the stable limit of the sum. Thus the already proved
positivity of every $a_{(t^{\rho'_j})}(h)$ implies
$a_{t\rho}(h)>0$.

Finally, $td\geq t|\rho|=|t\rho|$. Manivel stability, with the two
rectangle sides exchanged, therefore identifies
$g((t\rho)(htd),(td)^h,(td)^h)$ with $a_{t\rho}(h)$.
Since $(t\rho)(htd)=t\lambda$, this is (O-004.1).
\end{proof}

This concerns the stable-tail range $|\bar\lambda|\leq d$. The
padded-permanent condition only gives $|\bar\lambda|\leq md$, so this
is not a universal factor-two theorem for every GCT candidate.
The source's separate Corollary~5.4 asserts that for rectangle height
$h\geq7$, the group generated by the positive rectangular semigroup
contains every partition of size divisible by $h$ and length at most
$h^2$, the ambient representation-theoretic length bound. Together with its
full rational cone, this says its saturation is as large as possible;
it does not say the original semigroup is saturated. O-002 supplies
explicit holes in that maximal saturation.

\section{From formulas to genuine multiplicity statements}

Let $V=\mathbb C^{n^2}$, let
$\Omega_n=\overline{\mathrm{GL}(V)\cdot\det_n}$, and let
$Z_{n,m}=\overline{\mathrm{GL}(V)\cdot X_{11}^{n-m}\operatorname{per}_m}$.
The padding variable is internal: $X_{11}$ is one of the $m^2$
permanent variables, exactly as in both supplied IP and BIP sources.
An external padding variable would change the elementary length
bound to $m^2+1$ and is not silently substituted in their stated
thresholds. Put
\[
\begin{split}
 a_\lambda(d[n])&=\operatorname{mult}_{S_\lambda V}
                         \operatorname{Sym}^d(\operatorname{Sym}^nV),\\
 K_n(\lambda)&=\operatorname{mult}_{S_\lambda V}
                         \mathbb C[\Omega_n]_d,\\
 q^m_\lambda(d[n])&=\operatorname{mult}_{S_\lambda V}
                         \mathbb C[Z_{n,m}]_d.
\end{split}
\]
Here $|\lambda|=nd$.  The relevant chain is
\[
 K_n(\lambda)\leq\operatorname{sk}(\lambda,(d^n))
                  \leq g_{n,d}(\lambda),\qquad
 q^m_\lambda(d[n])\leq a_\lambda(d[n]).                    \tag{O-005.0}
\]
The middle term is the symmetric rectangular Kronecker coefficient,
the multiplicity in the coordinate ring of the determinant orbit;
it is not automatically the closure multiplicity. The chain and
interpretation are stated in the supplied B\"urgisser--Ikenmeyer--Panova
source, Section~1(b), and refer there to BLMW, equation~(5.2.7).

\subsection{O-005: exterior operations exclude every nontrivial hook}

The vanishing statement below is known: Giannelli,
\emph{On the decomposition of the Foulkes module}, Theorem~1.1;
the exterior-power identities are his Lemmas~3.2--3.3 and
Proposition~3.4. We record a short $\lambda$-ring form of its proof,
because this identifies exactly what the Burnside-to-representation
construction supplies here.

\paragraph{Finding O-005 (the $\lambda$-ring dictionary for hook vanishing).}
Let $P_N$ be the permutation representation of $S_N$ on $N=nd$
letters and $U_N=P_N-1$ its augmentation representation. If
$H=S_n\wr S_d$, then
\[
       \dim(\lambda^k U_N)^H=0\quad(1\leq k<N),
       \qquad \dim(\lambda^0 U_N)^H=1.                  \tag{O-005.1}
\]
By Schur--Weyl duality this says
\[
               a_{(nd-k,1^k)}(d[n])=0\quad(1\leq k<nd).
\]
Thus the nontrivial hook formulas for rectangular Kronecker coefficients
provide neither occurrence nor multiplicity obstructions for these
orbit closures: the permanent-side multiplicity already vanishes.

\begin{proof}[Proof of the $\lambda$-ring form]
Assume first $n,d\geq2$ and let $B=(S_n)^d\triangleleft H$.
As an $H$-module, $U_N$ is the direct sum of the block-constant
augmentation space $U_d$ and the direct sum $T$ of the $d$ within-block
augmentation spaces $U_n^{(i)}$.  The base group $B$ acts trivially on
$U_d$ and independently on the summands of $T$.

For $S_s$, the only invariant in the exterior algebra of $U_s$ has
degree zero. Indeed, $\Lambda^r P_s$ has no invariant for $r\geq2$:
a transposition of two members of an $r$-subset fixes that subset but
negates its exterior basis vector. For $r=1$ the invariant dimension
is one. The identity
$\lambda_t(P_s)=(1+t)\lambda_t(U_s)$ now proves the claim degree by
degree, using semisimplicity in characteristic zero.
The exterior-power direct-sum decomposition consequently gives
\[
       (\Lambda^\bullet U_N)^B=\Lambda^\bullet U_d.
\]
Taking invariants under $H/B=S_d$ gives (O-005.1).  If $n=1$ or
$d=1$, the wreath product is $S_N$ and the same claim applies directly.

Finally, $\lambda^k U_N=[nd-k,1^k]$, and the natural decomposition
\[
 \operatorname{Sym}^d\operatorname{Sym}^n V
       =(V^{\otimes nd})^{S_n\wr S_d}
       =\bigoplus_{\lambda\vdash nd}S_\lambda V
                         \otimes[\lambda]^H
\]
identifies its multiplicity with the left side of (O-005.1).
The quotient defining a coordinate ring cannot create new
characteristic-zero constituents, proving the obstruction conclusion.
\end{proof}

Two further known ambient restrictions are useful when testing
candidates. Giannelli's Proposition~2.7 gives
$a_\lambda(d[n])=0$ when $\ell(\lambda)>d$. His Theorem~1.2
also excludes many partitions near hooks: write
\[
 \lambda=(nd-k-|\alpha|,\alpha_1+1,\ldots,\alpha_t+1,1^{k-t}),
 \qquad t\leq k,\qquad b=\sum_{j=2}^{t}\alpha_j.
\]
For a nonempty partition $\alpha$, if $k>b$ and
$\alpha_1<(k-b)(k-b+1)/2$, then $a_\lambda(d[n])=0$.
These are known vanishing criteria, not new claims about saturation.

\subsection{O-006: full determinant multiplicities in two-row shapes}

\paragraph{Finding O-006.a (subspace criterion).}
Let $X\subseteq\operatorname{Sym}^n V^*$ be a closed
$\mathrm{GL}(V)$-stable subscheme over $\mathbb C$. Suppose it contains
the reduced linear space $\operatorname{Sym}^n U^*$ for some
$r$-dimensional quotient $V\twoheadrightarrow U$.
Then the coordinate-ring quotient loses no multiplicity with
$\ell(\lambda)\leq r$:
\[
 \operatorname{mult}_{S_\lambda V}\mathbb C[X]_d
                         =a_\lambda(d[n])
              \quad\text{if }\ell(\lambda)\leq r.       \tag{O-006.1}
\]

\begin{proof}
There is a decomposition natural in linear maps of vector spaces,
\[
 F(V):=\operatorname{Sym}^d\operatorname{Sym}^n V
       =\bigoplus_{\lambda\vdash nd}S_\lambda V\otimes M_\lambda,
 \qquad M_\lambda=[\lambda]^{S_n\wr S_d}.
\]
This follows by applying wreath-product invariants to Schur--Weyl
decomposition of $V^{\otimes nd}$; the factor $M_\lambda$ is independent
of $V$. The polynomial-functor decomposition and its naturality are
Macdonald, \emph{Symmetric Functions and Hall Polynomials}, second
edition, Chapter~I, Appendix~A, equations~(5.3)--(5.4), p.~155, and
(5.5)--(5.6), p.~156. Let $I(X)_d\subseteq F(V)$ be the degree-$d$ ideal.
Semisimplicity gives
$I(X)_d=\bigoplus_\lambda S_\lambda V\otimes J_\lambda$
for subspaces $J_\lambda\subseteq M_\lambda$.
Restriction to the displayed linear subspace is $F(\pi)$, where
$\pi:V\twoheadrightarrow U$. It acts on the $\lambda$ summand as
$S_\lambda(\pi)\otimes\operatorname{id}_{M_\lambda}$.
The map $S_\lambda(\pi)$ is surjective, since a linear section of
$\pi$ gives a right inverse after applying the Schur functor.
Its target is nonzero precisely when $\ell(\lambda)\leq r$.
Since every element of $I(X)_d$ restricts to zero, these facts force
$J_\lambda=0$ for such $\lambda$. This proves (O-006.1), including
its assertion about all copies of an irreducible, not just occurrence.
\end{proof}

\paragraph{Finding O-006.b.}
For $n\geq2$, $d\geq1$, and $\lambda\vdash nd$ with at most two
rows,
\[
 K_n(\lambda)=a_\lambda(d[n])
             =\operatorname{sk}(\lambda,(d^n))
             =g_{n,d}(\lambda).                         \tag{O-006.2}
\]
In particular, no such $\lambda$ can be a multiplicity obstruction
between the padded permanent and the determinant, for any $m<n$.

\begin{proof}
The polynomial map that substitutes an arbitrary linear endomorphism
of $V$ into $\det_n$ takes every endomorphism into $\Omega_n$:
$\mathrm{GL}(V)$ is dense in $\operatorname{End}(V)$, and a polynomial
that vanishes on its image also vanishes on the image of every
endomorphism. A binary homogeneous form of degree $n$ over
$\mathbb C$ factors into $n$ linear forms. It is the determinant of
the diagonal matrix with those forms on its diagonal, hence belongs
to $\Omega_n$. Consequently $\Omega_n$ contains every linear space
of binary forms, and (O-006.1) with $r=2$ proves $K_n=a_\lambda$.

For $\lambda=(nd-k,k)$ with $0\leq k\leq nd/2$, the weight
$(nd-k,k)$ in $\operatorname{Sym}^d\operatorname{Sym}^n\mathbb C^2$
has dimension $p_k(d,n)$: a monomial is a multiset of $d$ numbers
from $0,\ldots,n$ with sum $k$. The irreducible
$S_{(nd-j,j)}\mathbb C^2$ has each weight with second coordinate
$j,j+1,\ldots,nd-j$ once. Subtracting successive weight dimensions
therefore gives
\[
                  a_{(nd-k,k)}(d[n])=p_k(d,n)-p_{k-1}(d,n).
\]
The same difference is the rectangular Kronecker coefficient by
F-002 (the identity of Pak--Panova, Lemma~3). Thus $a_\lambda=g$.
The inequalities (O-005.0) imply the equality with the symmetric
coefficient as well. Finally $q^m_\lambda\leq a_\lambda=K_n$,
which excludes a multiplicity obstruction.
\end{proof}

\subsection{O-007: the same exclusion extends to three rows}

\paragraph{Finding O-007.}
For $n\geq2$ and $\ell(\lambda)\leq3$,
\[
                         K_n(\lambda)=a_\lambda(d[n]).   \tag{O-007.1}
\]
Thus a multiplicity obstruction for these two orbit closures must have
at least four rows. This strengthens the two-row exclusion; equality
with the ordinary rectangular Kronecker coefficient is not asserted
for three rows.

\begin{proof}
Beauville, \emph{Determinantal hypersurfaces}, Proposition~3.1(a),
printed p.~46, states that a smooth plane curve of degree $n$ with
a non-effective line bundle of degree $g-1$ has an $n\times n$
linear determinantal equation. Such a bundle exists over $\mathbb C$:
for $g\geq1$ the effective classes form the proper theta locus in
$\operatorname{Pic}^{g-1}(C)$; for $g=0$, take
$\mathcal O_{\mathbb P^1}(-1)$. The properness assertion also follows
from $\dim C^{(g-1)}=g-1<g=\dim\operatorname{Pic}^{g-1}(C)$.
For the relative Jacobian and theta locus see Beauville's following
paragraph~(3.2), on the same page.

Hence every smooth ternary form of degree $n$ is a determinant of
linear forms, up to a nonzero scalar that can be absorbed into a row
of the matrix. The substitution argument in O-006.b puts these
forms in $\Omega_n$. Smooth ternary forms are a dense open subset
of the space of all ternary forms (the discriminant is nonzero, as
the Fermat form verifies in characteristic zero). Since $\Omega_n$
is closed, it contains that entire space. Apply O-006.a with $r=3$.
\end{proof}

\subsection{O-008: full ambient multiplicity first fails in four rows}

\paragraph{Finding O-008.}
For every $n\geq4$, there exist $d\geq1$ and a partition
$\lambda\vdash nd$ with exactly four rows such that
\[
                         K_n(\lambda)<a_\lambda(d[n]).   \tag{O-008.1}
\]
Thus the universal equality in O-007 cannot be extended to every
four-row partition. This is an existence assertion about the
determinant ideal; it does not assert a permanent-versus-determinant
multiplicity obstruction.

\begin{proof}
Fix a four-dimensional quotient $V\twoheadrightarrow U$ and a linear
section. Let $L=\operatorname{Sym}^n U^*$, viewed as a linear subspace
of $\operatorname{Sym}^n V^*$, and let $r$ be the corresponding
linear retraction to $L$. Define
\[
 \delta_U:\operatorname{Mat}_n(\mathbb C)^4\longrightarrow L,
 \qquad (A_0,A_1,A_2,A_3)\longmapsto
                   \det(x_0A_0+x_1A_1+x_2A_2+x_3A_3),
\]
and put $D_U=\overline{\operatorname{im}\delta_U}$.
Every element of the image is an endomorphism substitution into
$\det_n$, so $D_U\subseteq\Omega_n\cap L$.
Conversely, restricting any endomorphism substitution to $L$ gives
a point of $\operatorname{im}\delta_U$. Since $r$ is continuous and
$\Omega_n$ is the closure of the endomorphism-substitution image,
$r(\Omega_n)\subseteq D_U$. A point of $\Omega_n\cap L$ is fixed by
$r$, proving
\[
                         \Omega_n\cap L=D_U.            \tag{O-008.2}
\]

We give an explicit parameter count for $D_U$, avoiding any assumption
that a dimension bound for an orbit quotient also controls its closure.
On the dense open set where $A_0$ is invertible, write
$B_i=A_0^{-1}A_i$ and $c=\det A_0$. On a further dense open set,
$B_1$ has distinct eigenvalues, so simultaneous conjugation makes it
diagonal. We can also assume that all entries $(B_2)_{1j}$ for
$j>1$ in an eigenbasis are nonzero. This condition is independent of
scaling eigenvectors, and holds on a nonempty open set after the
finite choice of ordering eigenvalues; for example take $B_1$
diagonal with distinct entries and $B_2$ with all entries nonzero.
Conjugation by a diagonal matrix then makes
$(B_2)_{1j}=1$ for $j>1$, while preserving the diagonal form of $B_1$.

Consequently a dense set of the image is parametrized by
\[
 c\det\bigl(x_0 I+x_1\operatorname{diag}(b_1,\ldots,b_n)
                                      +x_2 C+x_3 E\bigr),
 \qquad C_{1j}=1\ (j>1).
\]
The number of parameters is
$1+n+(n^2-n+1)+n^2=2n^2+2$.
The coefficients of these forms are polynomials in that many
parameters. Their generated field therefore has transcendence degree
at most $2n^2+2$, proving
$\dim D_U\leq2n^2+2$.
Equivalently, the original left-right action of
$\mathrm{SL}_n\times\mathrm{SL}_n$ has generic orbit dimension
$2n^2-2$: a tuple beginning with $I$, a diagonal matrix with distinct
entries, and a matrix with all entries nonzero has only a finite
simultaneous centralizer. The explicit parameter count already
suffices for the asserted bound.

For $n\geq4$,
\[
 \dim L-(2n^2+2)
 =\binom{n+3}{3}-2n^2-2
 =\frac{(n-1)(n-2)(n-3)}6>0.
\]
Thus $D_U$ is a proper closed cone. There is a nonzero homogeneous
polynomial $f\in\mathbb C[L]_d$ vanishing on it, with $d\geq1$.
The polynomial $f\circ r$ is nonzero and lies in $I(\Omega_n)_d$,
because $r(\Omega_n)\subseteq D_U$.
The injection
$\operatorname{Sym}^d\operatorname{Sym}^n U\hookrightarrow
\operatorname{Sym}^d\operatorname{Sym}^n V$ induced by the chosen
section has image only in isotypic components $S_\lambda V$ with
$\ell(\lambda)\leq4$, by the same functorial Schur decomposition as
O-006.a. At least one such component of $f\circ r$ is nonzero in the
ideal. O-007 excludes every component with at most three rows.
Therefore a component with exactly four rows occurs in the ideal,
which is precisely (O-008.1).
\end{proof}

\paragraph{Where the remaining multiplicity gate lies.}
The general sufficient inequality
$q^m_\lambda(d[n])>g_{n,d}(\lambda)$ implies a determinantal lower
bound (Ikenmeyer--Panova Theorem~1.4). It is stronger than the actual
multiplicity inequality $q^m_\lambda(d[n])>K_n(\lambda)$.
Ikenmeyer--Panova Proposition~2.8 already shows
\[
       q^m_\lambda(d[n])\leq a_\lambda(d[n])\leq g_{n,d}(\lambda)
                 \quad\text{if }n\geq|\bar\lambda|.
\]
Its Corollary~1.9 consequently requires $d>n/m$ for a
Kronecker-upper-bound multiplicity obstruction satisfying
$|\bar\lambda|\leq md$. This does not exclude an actual obstruction
caused by loss of multiplicity from the determinant orbit to its
closure. O-006 and O-007 do exclude that possibility in the specified
short-row families, because the closure has full ambient multiplicity.

\subsection{O-009: ambient-positive padded shapes already have positive $g$}

\paragraph{Finding O-009.}
Let $m\geq3$, $n>3m^4$, and $\lambda\vdash nd$. If
\[
 \ell(\lambda)\leq m^2,\qquad |\bar\lambda|\leq md,
                     \qquad a_\lambda(d[n])>0,
\]
then $g_{n,d}(\lambda)>0$.
This is a useful explicit variant of the known Ikenmeyer--Panova
no-occurrence theorem: membership in the actual permanent coordinate
ring is not required.

\begin{proof}
If $g_{n,d}(\lambda)=0$, then $a_\lambda(d[n])>g_{n,d}(\lambda)$.
Ikenmeyer--Panova Corollary~1.9 gives
$d>n/m>3m^3$. Their Theorem~1.10(a) excludes the exceptional tail set
\[
 X=\{(1),(1,1),(1,1,1,1),(1,1,1,1,1,1),(2,1),(3,1)\},
\]
because the ambient multiplicity is positive. All the hypotheses
of their Theorem~1.10(b) now hold, giving $g_{n,d}(\lambda)>0$,
a contradiction.
\end{proof}

Thus in this range adding ambient occurrence to the coarse padding
conditions eliminates the holes themselves. O-002 shows why that
additional hypothesis cannot be dropped. This statement does not
compare positive multiplicities on the two orbit closures, and it
does not prove ordinary saturation outside its stated region.

\subsection{O-010: a dimension bound on universal projection arguments}

\paragraph{Finding O-010.}
Let $r\geq3$, $n\geq m\geq1$, and fix a nonzero linear form $L$
on $\mathbb C^r$. If every form
$L^{n-m}f$, with $f$ an arbitrary homogeneous degree-$m$ form in
$r$ variables, is in the closure of the set of determinants of
$n\times n$ matrices of linear forms in those variables, then
\[
       \binom{m+r-1}{r-1}\leq(r-2)n^2+2.                \tag{O-010.1}
\]
Consequently, for a fixed positive integer $c$, a universal assertion
of this kind with $n=m^c$ and a fixed number $r$ of variables for
all sufficiently large $m$ requires $r\leq2c+1$. In particular
$r\leq51$ is necessary when $c=25$.

\begin{proof}
For $n\geq2$, apply the generic gauge from O-008 to an $r$-tuple
of matrices. After making its first matrix invertible and its
second diagonal, and normalizing $n-1$ entries of its third, the
remaining $r-3$ matrices are unrestricted. The number of parameters is
\[
       1+n+(n^2-n+1)+(r-3)n^2=(r-2)n^2+2.
\]
Thus the closed determinantal locus has dimension at most the
right side of (O-010.1). For $n=1$ the locus is the space of all
linear forms and has dimension $r=(r-2)n^2+2$, so the same bound
holds. Multiplication by the nonzero polynomial $L^{n-m}$ is an
injective linear map on the polynomial ring, so the space of forms
in the hypothesis is a linear space of dimension
$\binom{m+r-1}{r-1}$. Its containment proves (O-010.1).
For $n=m^c$, the left side is a polynomial in $m$ of degree $r-1$
with positive leading coefficient $1/(r-1)!$, while the right side
has degree $2c$. The inequality for all large $m$ therefore forces
$r-1\leq2c$.
\end{proof}

This bounds arguments that replace projected permanents by every
form of the same degree and number of variables. It does not bound
arguments using properties specific to the permanent, whose
projections occupy a much smaller parameter space. Combined with
the constructive bounds of this run, for $c\geq6$ it identifies
the optimal asymptotic fixed-variable threshold $r=2c+1$ for that
universal containment method. It is not a limitation theorem for
geometric complexity theory as a whole.

\subsection{O-011: an exact multiplicity-one family detected by Hessians}

\paragraph{Finding O-011.}
Let $n>m\geq2$, and, for $1\leq d\leq n^2$, put
\[
              \lambda_d=(d(n-2)+2,2^{d-1})\vdash nd.
\]
Then the actual closure and ambient multiplicities are
\[
 K_n(\lambda_d)=a_{\lambda_d}(d[n])=1,\qquad
 q^m_{\lambda_d}(d[n])=
 \begin{cases}1&d\leq m^2,\\0&m^2<d\leq n^2.\end{cases}       \tag{O-011.1}
\]
In particular this entire family, including outputs with as many
as $m^2$ rows and positive permanent multiplicity, has no actual
multiplicity obstruction. No bound such as $n\geq m^{25}$ or
the sufficient weight-box inequality is needed for this family.

\begin{proof}
For a form $F$ of degree $n$ on $V$, define the degree-$d$
polynomial function
\[
 W_d(F)=\det\left(\frac{\partial^2F}{\partial x_i\partial x_j}
                                  (e_1)\right)_{1\leq i,j\leq d}.
\]
For an upper triangular $g\in\mathrm{GL}(V)$, write $g_d$ for
its leading $d\times d$ block. Since $ge_1=g_{11}e_1$, the chain
rule and homogeneity give
\[
 W_d(g^{-1}\cdot F)
       =g_{11}^{d(n-2)}\det(g_d)^2 W_d(F).
\]
Here $g^{-1}\cdot F=F\circ g$, consistently with the action on
the coordinate ring. Hence $W_d$ is a highest-weight vector of
weight $\lambda_d$, provided it is nonzero. The ambient
multiplicity is one by F-011 with $k=2$: in the Cauchy
decomposition on $d$ columns, the relevant multiplicity module is
\[
 S_{\lambda_d}\mathbb C^d
       =\det^2\otimes\operatorname{Sym}^{d(n-2)}\mathbb C^d.
\]
Its column weight $(n,\ldots,n)$ is the one-dimensional space
spanned, after removing the determinant factor, by
$\prod_i z_i^{n-2}$, and column permutations fix it. The case
$d=1$ is the same calculation directly.

At the identity matrix, the Hessian of $\det_n$ is the bilinear form
\[
               (A,B)\longmapsto\operatorname{tr}(A)\operatorname{tr}(B)
                                      -\operatorname{tr}(AB).
\]
It is nondegenerate: the trace-zero subspace is nondegenerate for
$-\operatorname{tr}(AB)$, its orthogonal complement is the span
of $I_n$, and the latter has square $n(n-1)\ne0$.
A nondegenerate symmetric bilinear form and a vector with nonzero
square admit a nondegenerate $d$-dimensional subspace containing
that vector for every $1\leq d\leq n^2$; diagonalize its
orthogonal complement to see this. Choose coordinates taking
$e_1$ to $I_n$ and the first $d$ coordinate directions to such
a subspace. Then $W_d$ is nonzero on a translate of $\det_n$.
Thus $K_n(\lambda_d)\geq1$, and the ambient multiplicity-one
bound gives $K_n(\lambda_d)=1$.

It remains to prove nonvanishing on the internally padded permanent.
Set $s=n-m>0$, $P=\operatorname{per}_m$, and $F=X_{11}^sP$.
Work first at the all-ones matrix. Put $h=(m-2)!$, let $u$ be
the all-ones vector of length $m^2$, let $e$ be the coordinate
vector for the $(1,1)$ entry, and put
\[
 B=(J_m-I_m)\otimes(J_m-I_m).
\]
Direct differentiation of the permanent gives
$P=m!$, $\nabla P=(m-1)!u$, and $H_P=hB$ at that point.
Therefore
\[
 \frac{H_F}{h}
   =B+s(m-1)(eu^T+ue^T)+s(s-1)m(m-1)ee^T.              \tag{O-011.2}
\]
The eigenvalues of $B$ are $(m-1)^2$, $-(m-1)$ and $1$, with
multiplicities $1$, $2(m-1)$ and $(m-1)^2$, respectively, so it
is invertible. Its inverse gives the following $2\times2$ Gram matrix:
\[
 G=\begin{pmatrix}e^TB^{-1}e&e^TB^{-1}u\\
                  u^TB^{-1}e&u^TB^{-1}u\end{pmatrix}
   =\frac1{(m-1)^2}
       \begin{pmatrix}(m-2)^2&1\\1&m^2\end{pmatrix}.
\]
Apply the matrix determinant lemma to (O-011.2), with
\[
 C=\begin{pmatrix}s(s-1)m(m-1)&s(m-1)\\s(m-1)&0\end{pmatrix}.
\]
Expanding $\det(I_2+CG)$ and factoring yields the exact identity
\[
 \frac{\det H_F}{\det H_P}
  =\frac{m+s-1}{m-1}\bigl(1-(m^2-3m+1)s\bigr)
  =\frac{n-1}{m-1}\bigl(1-(m^2-3m+1)(n-m)\bigr).       \tag{O-011.3}
\]
For positive integral $s$ this is nonzero except when $(m,n)=(3,4)$:
the integer $m^2-3m+1$ is $-1$ for $m=2$, $1$ for $m=3$,
and at least $5$ for $m\geq4$.

The exceptional point does not cause an exceptional coefficient.
For $m=3,n=4$, evaluate instead at the matrix with $X_{11}=2$
and every other entry equal to one. Then $F=16$. In the
orthonormal basis of sums and differences under interchange of the
last two rows and last two columns, the Hessian has blocks
\[
 A=\begin{pmatrix}
 4&2\sqrt2&2\sqrt2&10\\
 2\sqrt2&0&4&2\sqrt2\\
 2\sqrt2&4&0&2\sqrt2\\
 10&2\sqrt2&2\sqrt2&4
 \end{pmatrix},\qquad
 C_0=\begin{pmatrix}0&-2\sqrt2\\-2\sqrt2&-4\end{pmatrix}
 \quad\text{twice},\qquad (4).
\]
For clarity the four invariant vectors for the first block are
$e_{11}$, $(e_{12}+e_{13})/\sqrt2$,
$(e_{21}+e_{31})/\sqrt2$, and
$(e_{22}+e_{23}+e_{32}+e_{33})/2$; replacing sums by the
row or column differences gives the two sign blocks and the
double-sign line. Direct differentiation gives the displayed
blocks. Their determinants are $576,-8,-8,4$, so
$\det H_F=147456\ne0$.

In every case we have found a point $v$ in the $m^2$-dimensional
permanent variable space with $F(v)\ne0$ and nondegenerate Hessian.
Euler's identity gives
\[
                          H_F(v,v)=n(n-1)F(v)\ne0.
\]
As for the determinant, choose any $d\leq m^2$ and a
nondegenerate $d$-subspace containing $v$, and then an invertible
change of ambient coordinates. This gives a translate of $F$
on which $W_d$ is nonzero. Hence $q^m_{\lambda_d}(d[n])\geq1$;
the ambient bound gives equality. For $d>m^2$, the length
$\ell(\lambda_d)=d$ violates the necessary internal-padding
length condition (IP Theorem~1.8), so $q^m_{\lambda_d}(d[n])=0$.
This proves every case of (O-011.1).
\end{proof}

The construction is the determinant of a restriction of the second
jet of the universal degree-$n$ section at a chosen point. It is
equivariant under the Borel preserving that point and flag; its
weight is the displayed determinant-line character. Its derivatives,
matrix entries and identities before the convenient orthonormal
change of basis are all defined over $\mathbb Z$. The multiplicity
interpretation uses the characteristic-zero Schur decomposition.
This places the explicit calculation in the equivariant locally
free-sheaf and Schur-functor dictionary of SGA~3, I, 4.7.2--4.7.3,
SGA~6, VI, Theorem~3.3, and cote~133, scan~028, without asserting
the same semisimple multiplicities in residue characteristics.

\section{Padding from a bounded degree preserves full multiplicity}

\subsection{O-012.a: a determinant bound obtained by splitting degree}

For $b\geq1$ and $r\geq2$, set
\[
 u=\left\lfloor\frac{b-1}{2}\right\rfloor,\qquad
 v=\left\lceil\frac{b-1}{2}\right\rceil,\qquad
 B(b,r)=\binom{r-1+u}{u}+\binom{r-1+v}{v}-1.
\]
Every degree-$b$ homogeneous form $f$ in $r$ variables, with a
chosen nonzero linear form $L$, admits a homogeneous linear
determinant of size $B(b,r)$ equal to
\[
                           L^{B(b,r)-b}f.               \tag{O-012.1}
\]
Consequently $L^{n-b}f\in\Omega_n$ whenever $n\geq B(b,r)$
and $r\leq n^2$. In one variable a diagonal determinant of size
$b$ suffices. These are integral identities in universal
coefficient variables.

\begin{proof}
Use $L$ and $z_1,\ldots,z_{r-1}$ as coordinates, and write
$f=\sum_{|\gamma|\leq b}c_\gamma z^\gamma L^{b-|\gamma|}$.
Make two copies of the monomial tree on the \emph{same} $r-1$
variables. The first has vertices of total degree at most $u$
and edges pointing away from its root; the second has vertices
of total degree at most $v$ and edges pointing toward its root.
The parent of a nonzero multi-index is obtained by subtracting
the last nonzero coordinate, so paths are unique. Label the
variable edges by $-z_j$, and identify the two roots into one
central vertex. Keep all other vertices distinct, even when
their multi-index labels are the same. There are $B=B(b,r)$ vertices.

Fix once and for all a rule that splits every multi-index of
degree at most $u+v=b-1$ as $\alpha+\beta$, with
$|\alpha|\leq u$, $|\beta|\leq v$; taking an initial segment
of its sorted variable word gives such a rule. For each term
with $|\gamma|\leq b-1$, insert the coefficient edge from
the prefix vertex $\alpha$ to the suffix vertex $\beta$ with
label $c_\gamma L$. For $|\gamma|=b$, remove one variable
$z_j$ by a fixed rule, split $\gamma-e_j=\alpha+\beta$,
and put label $c_\gamma z_j$ on that coefficient edge.
If several labels have the same endpoints, add them in the
matrix entry. All entries remain homogeneous linear forms.
Put an $L$-loop at every noncentral vertex; the central entry
is exactly the sum of its coefficient-loop labels.

As in M-007, every non-loop directed cycle uses the central
vertex and exactly one coefficient edge. Thus a determinant
term chooses one coefficient-edge summand and the unique
variable paths to and from it, with loops at all unused vertices.
For $|\gamma|=k<b$, the cycle has $k+1$ edges. Its sign
$(-1)^k$ cancels the sign on the $k$ variable edges, and its
contribution is $c_\gamma z^\gamma L^{B-k}$.
For $|\gamma|=b$, the cycle has $b$ edges, of which $b-1$
are tree edges and the last carries $c_\gamma z_j$.
The signs again cancel and the contribution is
$c_\gamma z^\gamma L^{B-b}$.
Summation gives (O-012.1). Since $r\geq2$, the two trees have
at least $u+1$ and $v+1$ vertices, so $B\geq u+v+1=b$;
all exponents are nonnegative. Adjoin the block $L I_{n-B}$
when a larger size $n$ is required.
\end{proof}

This complements the variable-splitting bound M-007. Splitting
degree is more efficient when the base degree is fixed and the
number of variables grows. One may use the smaller of the two
proved bounds. No assertion that either is an optimal general
determinantal representation is made.

\subsection{O-012.b: an arbitrary-rank padding inequality}

Let $U$ have dimension $r$, let $b\leq n$, and let
$\mu\vdash db$ have length at most $r$. Put
\[
                \lambda=\mu+d(n-b)e_1.
\]
If a closed $\mathrm{GL}(V)$-stable cone
$X\subseteq\operatorname{Sym}^n V^*$ contains every form
$x_1^{n-b}f$ with $f\in\operatorname{Sym}^b U^*$, then
\[
       \operatorname{mult}_{S_\lambda V}\mathbb C[X]_d
                                     \geq a_\mu(d[b]).  \tag{O-012.2}
\]
Here $U$ is the span of the first $r$ coordinate directions
in $V$. This is an arbitrary-rank form of the padding inequality
stated in Ikenmeyer--Panova Proposition~2.6(b), p.~6, citing
Kadish--Landsberg Theorem~1.3. The proof below makes the
normalization issue explicit.

\begin{proof}
Put $p=n-b$, and let
$T:\operatorname{Sym}^b U\to\operatorname{Sym}^n U$
be multiplication by $e_1^p$. For the upper Borel,
$gT=g_{11}^pTg$, so $\operatorname{Sym}^d T$ takes a
highest-weight vector of weight $\mu$ to one of weight
$\lambda$. It is injective.

Use the differential pairings
$\langle a,F\rangle=a(\partial)F|_0$ in each symmetric power.
The transpose $R$ of the padding map
$P:f\mapsto x_1^pf$ is $\partial_{e_1}^p$. On a monomial
$e^\alpha$ of degree $b$ one has
\[
                    RT(e^\alpha)
                    =\frac{(\alpha_1+p)!}{\alpha_1!}e^\alpha.
\]
Thus $RT$ is an invertible diagonal operator in characteristic
zero. Pullback to the padded linear space, after applying
$\operatorname{Sym}^dT$, is
$\operatorname{Sym}^d(RT)$, also invertible.
In particular, a basis of the degree-$d$ highest-weight space
of weight $\mu$ lifts to highest-weight vectors of weight
$\lambda$ whose restrictions to that padded space remain
linearly independent. Since that space lies in $X$, the lifted
vectors remain independent in $\mathbb C[X]_d$, proving
(O-012.2). A different conventional normalization of the
symmetric-power pairing changes the displayed diagonal factors
by nonzero constants and does not change this argument.
\end{proof}

Together O-012.a and O-012.b give
\[
             K_n(\mu+d(n-b)e_1)\geq a_\mu(d[b])
                 \quad\text{if }n\geq B(b,r).           \tag{O-012.3}
\]
The geometric statement concerns a whole padded linear space;
it does not assert identical evaluation of the lifted polynomial
on a single unchanged base form. Invertibility of the diagonal
pullback, rather than a false coefficient-normalization identity,
is what preserves all copies.

\subsection{O-012.c: stabilization after removing the last-row rectangle}

Let $d\geq2$ and $\lambda\vdash nd$ have at most $d$ rows;
pad it with zeros to length $d$. Define
\[
 k=\lambda_d,\qquad
 \rho=(\lambda_2-k,\ldots,\lambda_d-k),\qquad
 t=|\rho|,\qquad b=k+t.
\]
Zero parts of $\rho$ may be deleted. If $b\geq1$ and $n\geq b$,
put
\[
       \mu=(k+(d-1)t,\lambda_2,\ldots,\lambda_d)\vdash db.
\]
Then
\[
                         a_\lambda(d[n])=a_\mu(d[b]).  \tag{O-012.4}
\]
When $b=0$, the partition is the one-row shape and its ambient
multiplicity is one. For $\ell(\lambda)<d$ the definition gives
$k=0$, $b=|\bar\lambda|$; for $\ell(\lambda)=d$ it can
give a much smaller base degree than the total tail size.

\begin{proof}
Compute stable plethysm multiplicities using a $d$-dimensional
row space and a $d$-dimensional column space, with matrix
variables $Z=(z_{ij})$. By the Cauchy decomposition, the row
highest-weight space of weight $\lambda$ is the column module
$S_\lambda\mathbb C^d$. Its subspace of column multidegree
$(n,\ldots,n)$ fixed by column permutations has dimension
$a_\lambda(d[n])$.

Every polynomial in this row highest-weight space is divisible
by $\det(Z)^k$. To justify the assertion about every vector,
multiplication by that determinant power embeds the row
highest-weight space of weight $\lambda-k^d$ into the one
of weight $\lambda$. As column modules the map is the nonzero
map between the identical irreducibles
$\det^k\otimes S_{\lambda-k^d}\mathbb C^d$ and
$S_\lambda\mathbb C^d$, so it is an isomorphism.

After division, write $a=n-k$. The row weight is
$(da-t,\rho)$ and every column has degree $a$.
Each monomial has total degree $t$ outside the first row.
Consequently its exponent of $z_{1j}$ is at least $a-t$
in every column $j$. Since $a\geq t$, every such polynomial
is divisible by
\[
                         Q=\prod_{j=1}^d z_{1j}^{a-t}.
\]
This factor is fixed by the upper row unipotent group and by
column permutations. Division by it preserves the highest-weight
condition, subtracts $d(a-t)e_1$ from the row weight, and
reduces every column degree to $t$. Conversely multiplication
by $Q$ reverses the operation. Thus it is a bijection of the
two relevant column-permutation modules, including all their
copies and their signs.

Multiplying back by $\det(Z)^k$ gives row weight
$(k+(d-1)t,\lambda_2,\ldots,\lambda_d)=\mu$ and
column degree $b=k+t$. The determinant factor has the same
column sign character before and after, so the invariant
dimensions agree. This proves (O-012.4). The tuple $\mu$ is
a partition: $\rho_1\leq t\leq(d-1)t$, so its first
row is at least its second. Cauchy naturality then carries
the stable multiplicity to every actual vector-space dimension
in which $\lambda$ and $\mu$ fit.
\end{proof}

\subsection{O-012.d: full actual multiplicities and the four-copy family}

Under the hypotheses of O-012.c, set $r=\ell(\lambda)$.
If $r\leq n^2$ and
\[
                          n\geq\max\{b,B(b,r)\},       \tag{O-012.5}
\]
then
\[
                         K_n(\lambda)=a_\lambda(d[n]). \tag{O-012.6}
\]
For $r=1$ use the pure-power diagonal determinant instead of
$B(b,r)$. This excludes an actual multiplicity obstruction for
every padded-permanent closure, since $q\leq a=K_n$.

Indeed, $\lambda=\mu+d(n-b)e_1$, so O-012.3 gives
$K_n(\lambda)\geq a_\mu(d[b])=a_\lambda(d[n])$.
The reverse inequality is the ambient coordinate-ring quotient.

In particular, suppose $m\geq2$, $n\geq m^{25}$, and
$r\leq m^2$. Every shape satisfying $b\leq25$ in O-012.c
has full ambient determinant multiplicity. To verify the uniform
bound, use monotonicity in $b,r$:
\[
 B(b,r)\leq B(25,m^2)=2\binom{m^2+11}{12}-1.
\]
For $m=2$ the latter is $909<2^{25}$. For $m\geq3$,
\[
 2\binom{m^2+11}{12}
  \leq\frac{2(3m^2)^{12}}{12!}
  <m^{24}\leq m^{25},
\]
because $m^2+11\leq3m^2$ and $2\cdot3^{12}<12!$.
Also $n\geq m^{25}\geq25\geq b$.

For the explicit first higher-multiplicity example, take $d\geq6$,
$n\geq8$ and
\[
               \lambda=(d(n-2)-4,8,2^{d-2}).
\]
Here $k=2$, $t=6$, and $b=8$. Removing the column
$\det^2$ factor and using the binary weight count of F-013
(equivalently O-006.b) gives
\[
 a_\lambda(d[n])=p_6(d,n-2)-p_5(d,n-2)=11-7=4.
\]
Both rectangle bounds are inactive since $d,n-2\geq6$.
Therefore
\[
       n\geq\max\{8,B(8,d)\}\quad\Longrightarrow\quad
                         K_n(\lambda)=a_\lambda(d[n])=4. \tag{O-012.7}
\]
In the BIP range this holds for every $6\leq d\leq m^2$.
The statement settles all four determinant-side copies and
excludes $q>K_n$ there; it does not assert a value for the
permanent-side multiplicity between zero and four.

\paragraph{Geometric dictionary and scope.}
The determinant matrices define morphisms of affine coefficient
schemes over $\mathbb Z$. Multiplication by a chosen section
and contraction by its dual are maps of locally free symmetric
powers; the factorial diagonal calculation specifies their
characteristic-zero effect on all highest-weight copies.
The stabilization is a calculation in the coordinate algebra
of the matrix scheme: a determinant-line factor, followed by
division by a common first-row monomial, preserves the indicated
torus-weight and finite-group-invariant modules. The precise
background is the Cauchy identity, Macdonald I(4.3), p.~63,
and the functorial Schur decomposition I Appendix~A(5.3)--(5.6),
pp.~155--156, together with SGA~3 I4.7.2--4.7.3 and SGA~6
VI3.3. The related known padding statement is IP Proposition~2.6(b).
No claim is made that the factorial factors are invertible over
all of $\mathbb Z$. Integral polynomial identities and
characteristic-zero irreducible multiplicities remain separate.

The new gate is exact: a shape in this stabilized class can
only be an actual multiplicity obstruction if it fails
(O-012.5). In particular, at $n\geq m^{25}$ a surviving
maximal-row shape $\ell(\lambda)=d$ must have
$\lambda_d+\sum_{i=2}^d(\lambda_i-\lambda_d)>25$.
This is a full-multiplicity exclusion, including non-scalar
ambient multiplicities, rather than another occurrence test.
The later quadratic Chow bound M-010 excludes all maximal-row
shapes with $d\leq m^2$ throughout the BIP range as well.
Thus this particular BIP specialization is superseded by that
stronger degree result. O-012.a and O-012.b remain the inputs
to F-014 and F-015, which apply also at arbitrarily large outer
degree; the parameter-dependent statement (O-012.5) remains
valid in every stated range.

\subsection{O-013: the Hessian family at every outer degree}

Let $n\geq2$, $1\leq r\leq\min(d,n^2)$, and put
\[
          \lambda_{d,r}=(nd-2(r-1),2^{r-1})\vdash nd.
\]
Then
\[
           K_n(\lambda_{d,r})=a_{\lambda_{d,r}}(d[n])=1.
                                                        \tag{O-013.1}
\]
For internal padding and $n>m\geq2$, one moreover has
\[
 q^m_{\lambda_{d,r}}(d[n])=
 \begin{cases}1&r\leq m^2,\\0&r>m^2.\end{cases}          \tag{O-013.2}
\]
There is no upper bound on the outer degree $d$. Thus this
extends O-011 beyond the maximal-row case $d=r$ and gives
genuine positive-permanent examples even at degrees beyond
the independent quadratic Chow bound M-010.

\begin{proof}
The one-row case is the power of the first coefficient, so
assume $r\geq2$. Apply F-014 with second row $2$:
\[
 a_{\lambda_{d,r}}(d[n])=
 a_{\mu}(d[2]),\qquad
       \mu=(2(d-r+1),2^{r-1}).
                                                        \tag{O-013.3}
\]
We give the short multiplicity-free argument needed on the
right, rather than presuming a formula for all quadratic
plethysms. In the space of quadratic forms on $\mathbb C^r$,
the upper triangular group has a dense open orbit. Indeed,
on the open set where all leading principal minors of the
symmetric coefficient matrix are nonzero, symmetric Gaussian
elimination gives an upper-triangular congruence to a diagonal
matrix with nonzero entries. Over $\mathbb C$ a further
diagonal congruence changes that matrix to the identity.
Consequently the ratio of two nonzero polynomial highest-weight
functions of the same weight is a rational invariant of the
upper triangular group, and is constant on its open orbit.
The two functions are proportional. Every irreducible
multiplicity in the coordinate ring of quadratic forms is
therefore at most one.

For a degree-$n$ form $F$, let $W_r(F)$ be the Hessian minor
used in O-011, and define
\[
                  H_{d,r}(F)=W_r(F)F(e_1)^{d-r}.
                                                        \tag{O-013.4}
\]
This is a degree-$d$ polynomial function. The chain-rule
calculation in O-011 and the identity
$F(ge_1)=g_{11}^{n}F(e_1)$ for upper triangular $g$ show
that it is a highest-weight vector of weight
\[
  (r(n-2)+2+n(d-r),2^{r-1})=\lambda_{d,r}.
\]
For $n=2$, evaluation on a nonsingular diagonal quadratic
form with nonzero first coefficient makes (O-013.4) nonzero.
Thus the multiplicity at $\mu$ in (O-013.3) is exactly one,
and so is the ambient multiplicity at $\lambda_{d,r}$.

At the determinant point, O-011 supplies for every $r\leq n^2$
a change of coordinates under which $W_r$ is nonzero and
$F(e_1)$ is nonzero: the first direction corresponds to the
identity matrix, and the other directions span a nondegenerate
Hessian subspace containing it. Hence (O-013.4) is nonzero
on the determinant orbit, proving $K_n\geq1$. The ambient
upper bound gives (O-013.1).

For the internally padded permanent, O-011 supplies the same
two nonvanishing conditions for every $r\leq m^2$, including
the separately checked exceptional pair $(m,n)=(3,4)$.
Thus (O-013.4) is nonzero on that orbit as well, and
$q\geq1$. Again the ambient upper bound gives $q=1$.
If $r>m^2$, the permanent depends on only $m^2$ variables,
so the length restriction gives $q=0$.
\end{proof}

This is an exact scalar family at all outer degrees, with
an explicit surviving highest-weight function. It does not
claim that occurrence determines multiplicity when the ambient
coefficient exceeds one. The determinant-side proof uses
the actual orbit, so no replacement of closure multiplicities
by symmetric or ordinary rectangular coefficients occurs.

\paragraph{Geometric dictionary.}
Quadratic forms are points of the affine scheme of symmetric
bilinear matrices. Triangular congruence gives the displayed
open orbit after base change to $\mathbb C$; the resulting
character semi-invariants are torus-weight sections in the
grading/comodule correspondence of SGA~3, I, Proposition~4.7.3,
pp.~35--36. The Hessian and evaluation polynomials define
sections of the indicated Schur-type summands in the coordinate
algebra of forms. Their functorial decompositions are Macdonald,
I Appendix~A, (5.3)--(5.6), pp.~155--156, and the locally free
$K_0$ framework is SGA~6, VI, Theorem~3.3, pp.~392--393.
The scheme maps and derivative polynomials have integral
coefficients, but the Borel orbit argument, complete
reducibility, and Hessian nondegeneracy used here are
characteristic-zero assertions. Nonzero restriction to the
actual orbit means that the corresponding section survives
the coherent ideal quotient of its closure.

\section{A bounded-exponent determinant construction and a strict test case}

\subsection{O-014.a: degree splitting with exponent bounds}
Let $b\geq1$, $r\geq2$, and choose positive integer caps
\[
             c=(c_2,\ldots,c_r),\qquad c_2=b,
                    \quad c_i\leq b\quad(i\geq3).
\]
Write $u=\lfloor(b-1)/2\rfloor$, $v=\lceil(b-1)/2\rceil$,
and define
\[
 N_j(c)=\#\{\alpha\in\mathbb Z_{\geq0}^{r-1}:
                         \alpha\leq c,\ |\alpha|\leq j\},
 \qquad D(b,c)=N_u(c)+N_v(c)-1.                         \tag{O-014.1}
\]
Equivalently,
\[
 N_j(c)=[z^j]\frac{1}{1-z}\prod_{i=2}^r
                                      (1+z+\cdots+z^{c_i}).
\]
For every degree-$b$ form
\[
 f(L,z_2,\ldots,z_r)=
       \sum_{|\gamma|\leq b,\,\gamma\leq c}
                          a_\gamma z^\gamma L^{b-|\gamma|},
\]
there is a homogeneous linear matrix of size $D=D(b,c)$
whose determinant equals $L^{D-b}f$. Its entries are linear
forms in the form variables, with integral polynomial dependence
on the coefficient parameters.
In particular, every $L^{n-b}f$ in this bounded coefficient
space lies in $\Omega_n$ if $n\geq D$ and $r\leq n^2$.

\begin{proof}
Use the two trees of O-012.a, retaining only multi-indices
$\alpha\leq c$ in the first tree of height $u$ and the second
tree of height $v$. The retained sets are closed under taking
parents, so every root path remains present. They have $N_u(c)$
and $N_v(c)$ vertices; merge only their roots.
For a coefficient with $|\gamma|<b$, split its sorted word
into $\alpha+\beta=\gamma$ with $|\alpha|\leq u$ and
$|\beta|\leq v$, and insert the coefficient edge labeled
$a_\gamma L$. Both endpoints satisfy the caps because they
are coordinatewise bounded by $\gamma$. For $|\gamma|=b$,
remove one variable $z_i$, split $\gamma-e_i$ in the same
way, and label the coefficient edge by $a_\gamma z_i$.
Again both endpoints and every path are retained.

All tree edges carry the negative of their variable, all
noncentral unused vertices carry $L$-loops, and coincident
coefficient-edge labels are added. Every directed non-loop
cycle passes through the central vertex and exactly one
coefficient edge. The determinant signs cancel the tree-edge
signs just as in O-012.a, giving one term
$a_\gamma z^\gamma L^{D-|\gamma|}$ for each coefficient.
Their sum is $L^{D-b}f$. Since $c_2=b$, the first variable
alone supplies paths of full heights $u$ and $v$, so
$D\geq u+v+1=b$. No negative power of $L$ occurs.
Adding an $L I_{n-D}$ block gives the larger determinant.
\end{proof}

For a partition $\lambda$, use $b=\lambda_2$ and
$c=(\lambda_2,\ldots,\lambda_r)$. F-015 proves that this
whole restricted coefficient space is sufficient to preserve
every highest-weight copy, so that $n\geq D(b,c)$ implies
$K_n(\lambda)=a_\lambda(d[n])$. The argument uses the
second-row stabilization F-014 and the invertible diagonal
padding pullback of O-012.b. The size in (O-014.1) is at
most the unbounded size $B(b,r)$, and can be strictly smaller.

\subsection{O-014.b: a genuine positive-permanent example beyond the earlier gates}

Here is an exact test case demonstrating that the bounded
coefficient refinement can add an exclusion, including when
the ambient multiplicity is larger than one. Set
\[
 m=256,\quad r=m^2=65536,\quad b=32,\quad
 c=(32,2^{65534}),\quad n=D(32,c),\quad d=n.
\]
The exact integer is
\[
n=N_{15}(c)+N_{16}(c)-1.
                                                               \tag{O-014.2}
\]
Take the partition
\[
                 \lambda=(n^2-131100,32,2^{65534}).     \tag{O-014.3}
\]
Then
\[
              4\leq q^m_\lambda(d[n])\leq K_n(\lambda)
                     =a_\lambda(d[n]).
                                                               \tag{O-014.4}
\]
The comparison of every copy follows from F-015, not from
the explicit one-copy construction below.

To reproduce (O-014.2), put $R=65534$. Since $u=15$, $v=16$
and the first cap is inactive in these degrees,
\[
 N_j(c)=\sum_{k=0}^j(j-k+1)
          \sum_{t=0}^{\lfloor k/2\rfloor}
                         \binom Rt\binom{R-t}{k-2t}.
                                                               \tag{O-014.5}
\]
The inner sum selects the tail variables used twice and once;
the factor $j-k+1$ counts the possible first-coordinate
exponents in a monomial of total degree at most $j$.
Thus (O-014.5), followed by $n=N_{15}+N_{16}-1$, is a
finite integer certificate independent of decimal rounding.
The complete decimal integer is retained in
\texttt{boxed\_degree\_split\_certificate.json}.

\paragraph{Positive permanent multiplicity.}
For the internally padded permanent, choose the nondegenerate
Hessian point from O-011 and an orthogonal basis beginning
with its nonisotropic radial direction. Every leading Hessian
minor in that basis is nonzero, including $W_2$ and $W_r$;
the form itself is nonzero in the first direction. Therefore
\[
                    W_r(F)W_2(F)^{15}F(e_1)^{d-r-30}
                                                               \tag{O-014.6}
\]
is a nonzero degree-$d$ highest-weight function on the
permanent orbit. Its second row is $2+30=32$, its remaining
nonfirst rows are $2$, and its first row is determined by
total size $nd$. Its weight is exactly (O-014.3), proving
$q^m_\lambda>0$. The stronger bound $q^m_\lambda\geq4$
follows from M-012.3, whose hypotheses hold because $b=32$
and $d\geq r+30$. It retains four independent binary seed
classes after multiplication by one common nonzero Hessian
factor on the integral permanent coordinate ring.

For the ambient lower bound, use F-013's four-dimensional
highest-weight space at
\[
                    (r(n-2)-4,8,2^{r-2})
\]
in outer degree $r$. Multiplication by the nonzero polynomial
$W_2^{12}F(e_1)^{d-r-24}$ sends all four vectors to the
highest-weight space at (O-014.3). The ambient polynomial
ring is a domain, so multiplication by this fixed polynomial
is injective. Thus $a_\lambda\geq4$. No value for the actual
exact permanent multiplicity is inferred from this ambient calculation.

\paragraph{Comparison with previously proved sufficient conditions.}
Exact integer arithmetic in the accompanying certificate gives
\[
                   m^{25}\leq n<2^m-1,
\]
so the BIP range holds and the full Grenet containment gate
F-009 does not apply. The unbounded size is
\[
B(32,65536)=\binom{65550}{15}+\binom{65551}{16}-1>n,
\]
so the earlier second-row determinant criterion also fails.
For any split of the $r-1=65535$ tail variables, some part
has at least $32768$ variables. Since all caps are at least
two, its monomial count through degree $m$ is at least
$\binom{32768}{256}$, by choosing squarefree monomials of
degree $m$. Consequently
\[
        T_m(\lambda)\geq\binom{32768}{256}>n.
\]
The analogous degree-$b$ variable-splitting bound is at least
$\binom{32768}{32}>n$. The uniform row cutoff does not apply,
since $r>51$, and the second row is greater than $25$.
Since $d=n$, the quadratic Chow bound $n\geq d(d-1)$
and the adaptive bound $n\geq B_d(32)=32d-497$ also fail.
The coarse padding requirements hold: $r=m^2$ and
$|\bar\lambda|=131100\leq md$.

These failures describe the scope of the comparison: the
example adds to the explicitly listed sufficient criteria.
It does not prove that every other conceivable method would
fail. The conclusion is the exact full determinant multiplicity
in (O-014.4), with at least four actual permanent copies.
The exact permanent multiplicity remains undetermined.
The example certifies strict refinement of the listed explicit
conditions; it is not an exhibited surviving $q>K_n$ candidate.
The later nested-support criteria F-016 and M-011 also handle
this example. O-015.c gives a separate four-copy example
beyond all the stabilization criteria, using M-011 directly.

\paragraph{Geometric dictionary.}
The bounded coefficient space is an affine linear subscheme
of the scheme of degree-$b$ forms, cut out by coefficient
coordinates outside the prescribed exponent box. Its coordinates
are selected by torus weights, using SGA~3, I, Proposition~4.7.3,
pp.~35--36. The two trees define a morphism from that affine
coefficient scheme to the scheme of linear matrices over
$\mathbb Z$, and the determinant identity proves factorization
through the padded-form morphism. F-015 then interprets
restriction as a torus-equivariant map of coordinate modules;
representations as comodules are SGA~3, I, Proposition~4.7.2,
pp.~34--35, and Schur functor naturality is Macdonald,
I Appendix~A, (5.3)--(5.6), pp.~155--156. The integral
factorization does not itself split irreducibles integrally:
the full-copy comparison is made after base change to
$\mathbb C$, with the explicit factorial invertibility in
O-012.b. Multiplication of the Hessian sections takes place
in the actual ambient coordinate algebra, and the positivity
claim comes from their nonzero restriction to the permanent
orbit rather than a formal identity in its $K_0$ class.

\section{Determinant graphs for nested highest-weight supports}

\subsection{O-015.a: arbitrary downward-closed coefficient supports}
Fix a base degree $b\geq1$ and $r\geq2$. Let
$\Gamma\subseteq\mathbb Z_{\geq0}^{r-1}$ be a nonempty finite
coordinatewise downward-closed set contained in $|\gamma|\leq b$.
Consider all forms
\[
             f(L,z)=\sum_{\gamma\in\Gamma}
                         a_\gamma z^\gamma L^{b-|\gamma|}.
\]
Put $u=\lfloor(b-1)/2\rfloor$, $v=\lceil(b-1)/2\rceil$,
$\Gamma_j=\{\gamma\in\Gamma:|\gamma|\leq j\}$, and
\[
                    D_\Gamma=|\Gamma_u|+|\Gamma_v|-1.
                                                               \tag{O-015.1}
\]
There is a homogeneous linear matrix of size $D_\Gamma$
whose determinant is the polynomial
\[
                 \sum_{\gamma\in\Gamma}
                         a_\gamma z^\gamma L^{D_\Gamma-|\gamma|}.
                                                               \tag{O-015.2}
\]
For a split $A\sqcup B=\{2,\ldots,r\}$, let $\Gamma_A$
and $\Gamma_B$ be the coordinate projections of $\Gamma$,
including the zero index. There is also such a matrix of size
\[
                         S_{\Gamma,A}=|\Gamma_A|+|\Gamma_B|-1.
                                                               \tag{O-015.3}
\]
In (O-015.2) replace $D_\Gamma$ by this size. In each case
all displayed exponents are nonnegative. Thus every
$L^{n-b}f$ in the prescribed coefficient space belongs to
$\Omega_n$ when $n\geq b$ and $n$ is at least one of the
displayed sizes, provided $r\leq n^2$.

\begin{proof}
For (O-015.1), use the two degree-split trees on
$\Gamma_u$ and $\Gamma_v$. Deleting a variable from a
multi-index stays in $\Gamma$, so all root paths exist.
Every split $\gamma=\alpha+\beta$ of a retained monomial
also has $\alpha,\beta\in\Gamma$, since each is coordinatewise
bounded by $\gamma$. For $|\gamma|<b$, use the coefficient
edge $a_\gamma L$; for $|\gamma|=b$, first remove one
variable $z_i$, split the remaining degree $b-1$, and use
$a_\gamma z_i$. These are precisely the cycles of O-014.a.
When $|\gamma|<b$ the cycle has $|\gamma|+1$ vertices;
when $|\gamma|=b$ it has $b$ vertices. Since these vertices
are distinct in the constructed graph, the corresponding
exponent in (O-015.2) is nonnegative even if the total size
is less than $b$. Permutation signs cancel all tree-edge
signs and yield (O-015.2).

For (O-015.3), the projections are downward closed, and
the projections of each $\gamma\in\Gamma$ give the two
required vertices. Use the variable-split trees on these
two sets and a coefficient edge $a_\gamma L$ between them.
Each cycle has $|\gamma|+1$ vertices, so its contribution
is $a_\gamma z^\gamma L^{S_{\Gamma,A}-|\gamma|}$, again
with nonnegative exponent. Padding either matrix by an
$L$-diagonal block of the required size gives
$L^{n-b}f$. Formally the unpadded identity can be written
$L^b\det M=L^{\operatorname{size}(M)}f$; the termwise
argument proves polynomiality even when $\operatorname{size}(M)<b$.
All identities are integral in the coefficient parameters.
\end{proof}

\subsection{O-015.b: nested support counts}
For a partition $\lambda$ with $r\geq2$ rows and
$b=\lambda_2$, take
\[
 \Gamma(\lambda)=\left\{\gamma\in\mathbb Z_{\geq0}^{r-1}:
              \sum_{i=j}^r\gamma_i\leq\lambda_j
                         \quad(2\leq j\leq r)\right\}.
                                                               \tag{O-015.4}
\]
The $j=2$ inequality bounds the total degree by $b$;
all inequalities are preserved by coordinatewise decrease.
Thus O-015.a applies. The pure second-variable powers through
degree $b$ belong to this set. Consequently $D_\Gamma\geq b$
and every $S_{\Gamma,A}\geq b+1$.
F-016 proves that this support contains the whole stabilized
highest-weight space, so either size bound preserves all
ambient determinant multiplicities at the corresponding
weight. The suffix conditions strengthen coordinatewise caps.

For $r\geq3$ write $c=\lambda_3$ and $R=r-2$.
For $0\leq j\leq b$, ignoring the later suffix inequalities
gives the useful bound
\[
 |\Gamma_j(\lambda)|\leq
 \sum_{h=0}^{\min(c,j)}(j-h+1)\binom{R+h-1}{h}.
                                                               \tag{O-015.5}
\]
Here $h$ is the total degree in variables $3,\ldots,r$;
the second variable then has $j-h+1$ possible exponents.
Equality holds if $\lambda_3=\cdots=\lambda_r=c$.
At the variable split $A=\{2\}$, $B=\{3,\ldots,r\}$,
\[
                   S_{\Gamma,A}\leq
                          b+\binom{r-2+c}{c}.          \tag{O-015.6}
\]
Again equality holds for the constant later tail. In particular,
for tails $(b,2^{r-2})$ the sufficient full-multiplicity bound
from F-016 is simply
\[
                       n\geq b+\binom r2.             \tag{O-015.7}
\]
For comparison the degree-split count in this case, for $2\leq j\leq b$,
is
\[
              |\Gamma_j|=(j+1)+j(r-2)+(j-1)\binom{r-1}{2}.
\]
The variable split can be substantially smaller, so both
constructions are retained.

\subsection{O-015.c: an actual four-copy exclusion beyond inner-degree stabilization}

Set
\[
 m=256,\quad r=65536,\quad n=m^{25}=2^{200},\quad
 d=3n,\quad b=2n,
\]
and take
\[
           \lambda=(3n^2-2n-131068,2n,2^{65534}).       \tag{O-015.8}
\]
Then the actual multiplicities satisfy
\[
                         4\leq q^m_\lambda(d[n])
                                      \leq K_n(\lambda).        \tag{O-015.9}
\]
Neither exact multiplicity is claimed here.

For the upper comparison, use M-011 at the cut $k=3$.
Although $\lambda_2=2n>n$, the retained permanent-degree
forms have total degree at most two in variables $3,\ldots,r$.
The prefix variable has degree at most $m$, so the required
matrix size is
\[
               m+\binom r2=2147451136<2^{200}=n.
\]
This proves $q\leq K_n$ without any inner-degree stabilization.
For the lower bound, M-012.3 applies with the same even $b$:
$d\geq r+b-2$ because $n\geq r-2$. It gives four
independent highest-weight classes on the permanent closure,
proving $q\geq4$. More explicitly, multiply the four binary
seed classes by
\[
              W_r\,W_2^{n-4}F(e_1)^{n-r+2}.
\]
The common factor is nonzero on the reduced irreducible
permanent closure, and the degree and weight are those in
(O-015.8). The seed classes each have degree six; the
displayed factor has degree $d-6$.

This example satisfies the BIP threshold with equality, while
$n<2^{256}-1$, so the full Grenet containment criterion does
not apply. All the stabilization-based full-multiplicity
criteria F-014--F-016 require $n\geq\lambda_2$ and fail.
For every variable split in the older M-009 box criterion,
one part has at least $32768$ variables, all with caps at
least two. Therefore
\[
                 T_m(\lambda)\geq\binom{32768}{256}>2^{200}.
\]
The quadratic Chow criterion fails since $d(d-1)>n$.
The adaptive one fails since, with $b<d$,
\[
                    B_d(b)=db-\binom b2-1
                                =4n^2+n-1>n.
\]
The coarse GCT shape conditions hold: $r=m^2$ and
$|\bar\lambda|=2n+131068\leq md$. In particular the
necessary positive-multiplicity condition $q\geq2$ is
verified rather than left unevaluated.

The accompanying exact integer certificate checks every
comparison. This certifies strict refinement of the listed
explicit sufficient conditions by M-011. It exhibits an
excluded shape with genuine multiplicity at least four; it
does not exhibit an actual obstruction or establish a value
for either of its two multiplicities.

\paragraph{Geometric dictionary.}
The downward-closed support defines an affine coefficient
subscheme selected by torus weights, in the grading/comodule
dictionary of SGA~3, I, Proposition~4.7.3, pp.~35--36.
The graph morphisms to linear matrices and their determinant
identities are defined over $\mathbb Z$. The supports need
not be stable under the entire general linear group: the
highest-weight and ideal arguments are performed in the
ambient representation, using SGA~3, I, Proposition~4.7.2,
pp.~34--35, and functorial Schur decomposition as in
Macdonald, I Appendix~A, (5.3)--(5.6), pp.~155--156.
The characteristic-zero all-copy restrictions are F-016 and
M-011. The four surviving permanent sections in the example
come from the explicit integral binary substitution of M-012,
followed by multiplication in the integral coordinate algebra
of the reduced irreducible complex orbit closure.

\section{A boundary for exact universal padded determinant constructions}

\subsection{O-016: the degree-split exponent is optimal for this exact method}
Fix $b\geq2$, and let $R\geq5$ be the number of variables.
Write
\[
 h=\lfloor b/2\rfloor,\qquad
 D_j(R)=\binom{R+j-1}{j}.
\]
For a Zariski-general homogeneous degree-$b$ form $f$ in
$R$ variables over $\mathbb C$, its exact affine-linear
determinantal complexity satisfies
\[
 \operatorname{dc}(f)\geq
 1+\left\lceil\frac{D_b(R)}{D_h(R)+D_{b-h}(R)}\right\rceil
                         =\Omega_b(R^{\lfloor b/2\rfloor}).
                                                               \tag{O-016.1}
\]
Consequently any construction representing every padded
degree-$b$ form in $r=R+1$ variables by an exact homogeneous
linear determinant needs size
$\Omega_b(r^{\lfloor b/2\rfloor})$. The size
$B(b,r)=O_b(r^{\lfloor b/2\rfloor})$ in O-012.a therefore
has the optimal exponent in $r$ for fixed $b$, within this
exact universal problem.

The known input is Chatterjee--Kumar--Volk,
\emph{Determinants vs. Algebraic Branching Programs}, ITCS~2024,
article~27: Observation~2.6, p.~27:6, and Theorem~3.1,
p.~27:8. They imply that if a homogeneous form has singular
locus of dimension less than $R-4$ and an exact determinant
representation of size $s$, it has a homogeneous algebraic
branching program of width at most $s-1$. Their Definition~2.7,
p.~27:7, counts vertices rather than edges for program size;
Theorem~4.1, p.~27:9, more generally gives size $O(b^5s)$
for all homogeneous forms. The sharper width statement is
the only complexity input needed below. These are known
theorems, not new results of this investigation.

\begin{proof}
There is a nonempty Zariski-open set of degree-$b$ forms
whose projective hypersurface is smooth. To see both openness
and nonemptiness, the incidence defined by all partial
derivatives in the product of the coefficient space with
$\mathbb P^{R-1}$ has closed projection, since the second
factor is projective (EGA~II, Theorem~5.5.3(i), p.~104,
with properness Definition~5.4.1, p.~100), while the Fermat form
$z_1^b+\cdots+z_R^b$ is outside that projection.
For such a form the affine singular locus is exactly the
origin, of dimension zero, which is less than $R-4$.

By the cited regular-representation result, any determinant
representation of such an $f$ of size $s$ gives a homogeneous
branching program of width at most $s-1$. Cut the program
at layer $h$. Its path sum then has the form
\[
                f=\sum_{i=1}^{s-1}P_iQ_i,
       \qquad \deg P_i=h,\quad\deg Q_i=b-h,             \tag{O-016.2}
\]
inserting zero summands if the actual width is smaller.
This is simply the partition of all source-to-sink paths by
their vertex at that layer.

For a fixed integer $s$, the forms in (O-016.2) lie in the
image of the polynomial morphism
\[
 \left(\operatorname{Sym}^h\mathbb C^R
             \times\operatorname{Sym}^{b-h}\mathbb C^R\right)^{s-1}
 \longrightarrow\operatorname{Sym}^b\mathbb C^R,
                   \quad (P_i,Q_i)_i\longmapsto\sum_iP_iQ_i.
\]
The closure of its image has dimension at most
$(s-1)(D_h(R)+D_{b-h}(R))$. If this is strictly smaller
than $D_b(R)$, that image closure is proper. Its union with
the proper closed nonsmooth locus is still proper, since
the ambient affine coefficient space is irreducible.
Outside the union no form has determinant complexity at most
$s$. Taking the largest integer $s$ satisfying the strict
dimension inequality proves (O-016.1).

For fixed $b$, the numerator has degree $b$ in $R$ and
the denominator has degree $\lceil b/2\rceil$, with
positive leading coefficients. This gives the stated
asymptotic exponent. Finally, an exact identity
\[
                    \det M(L,z)=L^{s-b}f(z)
\]
with $M$ homogeneous linear and $f$ independent of $L$
specializes at $L=1$ to an affine-linear determinant for
$f(z)$ of size $s$. A universal padded construction includes
these forms, so it inherits the lower bound. O-012.a has
the matching exponent because
$\lceil(b-1)/2\rceil=\lfloor b/2\rfloor$.
\end{proof}

\paragraph{The base-degree twenty-six threshold.}
Take $r=m^2$ and $b=26$, so $R=m^2-1$. Then
\[
 \frac{D_{26}(R)}{2D_{13}(R)}
   =\frac{13!}{2\cdot26!}\prod_{j=13}^{25}(R+j)
                         =\Theta(m^{26}).              \tag{O-016.3}
\]
For all sufficiently large $m$, this exceeds $m^{25}$.
Thus not every degree-$26$ form in these variables can be
handled by an exact padded determinant of size $m^{25}$.
The uniform second-row constant twenty-five obtained from
an exact universal base-form construction cannot be raised
to twenty-six by merely improving the constant factor in
that universal exact construction. Restricted supports and
permanent-specific families remain available, as already
used in F-015--F-017 and M-011.

\paragraph{Exact complexity versus orbit-closure containment.}
The last conclusion is deliberately about exact determinant
identities. It does not prove that the corresponding padded
linear space is absent from $\Omega_{m^{25}}$, nor that a
multiplicity obstruction exists. A point of an orbit closure
can be a limit of determinants whose specialization at $L=1$
has higher-degree terms that disappear only in the limit.
The homogeneous fixed-degree conversion theorem does not
automatically apply uniformly to those approximating forms.
The source itself discusses the unresolved size-preserving
homogeneous-extraction issue on p.~27:5. No passage from
exact to border determinantal complexity is assumed here.
The generic lower bound is also not an explicit lower bound
for the permanent or for any identified GCT coefficient.

\paragraph{Geometric dictionary.}
The coefficient and product parameter spaces above are affine
schemes, and the dimension argument bounds the dimension of
a scheme-theoretic image after base change to $\mathbb C$.
The smooth-form open set is the complement of the proper
projection of a projective incidence scheme. The padded
matrix identity is an integral morphism, while specialization
at the nonzero section $L=1$ is a fiber operation.
These geometric operations do not identify an image with
its closure or assert lifting through a limiting family.
The distinction is the same coherent-quotient versus explicit
parametrization distinction used throughout the GCT comparison.

\section{A border-valid dimension boundary for the mixed construction}

\subsection{O-017: the mixed staircase reaches the universal cut-family exponent}
Fix integers $k\geq2$, $s\geq1$, and a real exponent
$c\geq1$. Put $r=m^2$, and suppose $n=n(m)\geq m$
with $n(m)=O(m^c)$. Take one padding variable $L$,
$k-2$ early variables, and $R=r-k+1$ later variables.
For sufficiently large $m$, consider the whole linear space
of degree-$m$ forms having total later-variable degree at
most $s$. Its dimension is
\[
 D_m(k,s)=\sum_{h=0}^s
       \binom{R+h-1}{h}\binom{m-h+k-2}{k-2}
        =\frac{m^{2s+k-2}}{s!(k-2)!}
                                +O_{k,s}(m^{2s+k-3}).    \tag{O-017.1}
\]
If
\[
                            k+2s-4>2c,                 \tag{O-017.2}
\]
then not every form in that space, after multiplication by
$L^{n-m}$, belongs to $\Omega_n$ for all sufficiently large
$m$. This is a dimension statement about the actual orbit
closure, and therefore allows border representations.

\begin{proof}
For later degree $h$, the first binomial coefficient in
(O-017.1) counts the later monomials. The remaining degree
$m-h$ is distributed among the $k-2$ early variables and
$L$, giving the second binomial coefficient. Distinct
monomials are linearly independent, so the sum is the exact
dimension. Here $m\geq s$ is harmless because the conclusion
is asymptotic with fixed $s$. Since $R=m^2-k+1$,
the term $h=s$ has the displayed positive leading coefficient
and degree $2s+k-2$ in $m$; all smaller $h$ have lower
degree. Multiplication by the fixed nonzero monomial
$L^{n-m}$ is injective, so it preserves that dimension.

O-010 gives the dimension bound
\[
 \dim\bigl(\Omega_n\cap\operatorname{Sym}^n((\mathbb C^r)^*)\bigr)
                              \leq (r-2)n^2+2.         \tag{O-017.3}
\]
Its use here does not replace the orbit closure by the exact
image of substitutions. Indeed the intersection is the
\emph{closure} of the $r$-variable determinant-pencil image:
the projection to the first $r$ variables sends the full
determinant substitution image to that pencil image, and
continuity sends its closure into the closure of the latter.
A point of the intersection is fixed by this projection.
The reverse inclusion follows because each pencil is an
endomorphism substitution of the determinant. O-010 bounds
the dimension of this image closure by its explicit generic
gauge parametrization.

If the whole padded linear family were contained in
$\Omega_n$, its dimension would be at most the right-hand
side of (O-017.3). Since $r=m^2$ and $n=O(m^c)$, that
right-hand side is $O(m^{2c+2})$. Condition (O-017.2)
is exactly $2s+k-2>2c+2$. The positive leading term in
(O-017.1) then contradicts the required dimension inequality
for all sufficiently large $m$.
\end{proof}

\paragraph{Meaning for the staircase.}
The BIP staircase in F-017 uses $c=25$ and
$k=54-2s$, so $k+2s-4=50=2c$. The IP staircase
uses $c=4$, $k=12-2s$, again giving equality; a fixed
factor such as $n=3m^4+1$ does not change the exponent.
Thus those constructions reach the allowed exponent for
their universal cut families. Increasing the permitted later
degree by one at the same cut crosses (O-017.2).
No construction proving containment of that \emph{whole}
larger truncated family in the determinant orbit closure can
work at the same power range for all large $m$.

For a finite illustration in the IP range, take $m=100$,
$r=10000$, $n=3m^4+1$, $k=2$, and the raised cap $s=6$.
Then
\[
 D_m(2,6)=\binom{10005}{6}>(9998)(300000001)^2+2.
\]
The exact comparison is in the accompanying certificate.
For a BIP illustration, take $m=10^{15}$, $n=m^{25}$,
$k=4$, and the raised cap $s=26$; the same certificate
checks $D_m(4,26)>(m^2-2)m^{50}+2$ by integer arithmetic.

This boundary does not apply to a smaller family selected
by additional nested row bounds, nor to the actual family
of projected permanents. The dimension of an unrestricted
coefficient family is not the dimension or evaluation rank
of those smaller families. In particular, (O-017.2) does
not exhibit an actual multiplicity obstruction, and it does
not say that the F-017 staircase is a complete classification.
It identifies the limitation of extending the same universal
cut-family containment argument by one unit in the suffix cap.

\paragraph{Geometric dictionary.}
The truncated family is a finite torus-weight affine linear
subscheme, with the grading interpretation of SGA~3, I,
Proposition~4.7.3, pp.~35--36. Padding by a fixed section
is a closed linear embedding. The determinant side is a
scheme-theoretic image closure after base change to
$\mathbb C$, and its dimension is bounded by the explicit
parametrization in O-010. This comparison already includes
limits and makes no normality, cohomological purity, or
saturation assumption. It is separate from O-016's stronger
fixed-base-degree exponent statement, whose method requires
exact determinant representations.

\section{Using two later-row bounds simultaneously}

\subsection{O-018.a: a joint-row gate}
In the BIP numerical range $n\geq m^{25}$, every actual
multiplicity obstruction satisfies
\[
                         \lambda_3\geq139
                           \quad\hbox{or}\quad\lambda_{52}\geq3.
                                                               \tag{O-018.1}
\]
The numerical constant is a sufficient bound obtained from
the count below; it is not claimed optimal.
A-025 subsequently strengthens $139$ to $557$ by using
two support trees and retaining the total-degree constraint.
The elementary count here remains useful for the explicit
full-profile example below.

\begin{proof}
The earlier reductions already give $m\geq190$ and
$52\leq r=\ell(\lambda)\leq m^2$ for an actual obstruction.
Suppose instead that $\lambda_3\leq138$ and
$\lambda_{52}\leq2$. In a projected degree-$m$ form,
the nested support F-016 relevant to a highest-weight
function has total degree at most $138$ in variables
$3,\ldots,r$, and at most two in variables $52,\ldots,r$.
There is one earlier variable, $x_2$, forty-nine middle
variables, and $r-51$ last variables.

Let $\Gamma$ be the set of all nonfirst-variable multi-indices
of degree at most $m$ satisfying the full nested bounds.
It is nonempty and downward closed. Ignoring coupled degree
conditions only increases its size, so
\[
 |\Gamma|\leq(m+1)\binom{187}{49}\binom{r-49}{2}
                           \leq\binom{187}{49}m^5.     \tag{O-018.2}
\]
The second inequality uses $m+1\leq2m$ and
$\binom{r-49}{2}\leq r^2/2\leq m^4/2$.
Exact integer arithmetic gives
\[
                         \binom{187}{49}\leq190^{20}.
\]
Since $m\geq190$, (O-018.2) is at most $m^{25}\leq n$.

Place all of $\Gamma$ in one monomial tree, and take the
other tree to consist of its root. O-015.a gives a determinant
of size $|\Gamma|$ for every retained padded form.
For every determinant-ideal highest-weight vector, first align
the padding with $x_1$ as in M-009, then delete the coefficient
monomials outside $\Gamma$. F-016 preserves its evaluation,
and the retained form is the determinant just constructed.
Thus every such vector vanishes on the permanent closure.
The inclusion of all highest-weight ideal multiplicity spaces
gives $q\leq K_n$, a contradiction. This proves (O-018.1).
This is also a special case of the general support-factorization
criterion F-018, with one factor equal to the whole support
and the other consisting only of zero.
\end{proof}

\subsection{O-018.b: a four-copy example beyond every single-cut containment test}
Set
\[
 m=2^{200},\quad r=m^2,\quad n=m^{25}=2^{5000},\quad d=3n,
\]
and define
\[
 \lambda=(3n^2-2n-2r-1172,\ 2n,\ 26^{49},\ 2^{r-51}).
                                                               \tag{O-018.3}
\]
It has $r$ rows and size $nd$. Then
\[
                              4\leq q^m_\lambda(d[n])
                                              \leq K_n(\lambda).
                                                               \tag{O-018.4}
\]
The exact values of these multiplicities are not asserted.
Every single-cut universal coefficient family associated to
this partition has dimension greater than the corresponding
determinant-pencil closure. Thus the comparison in (O-018.4)
uses information beyond every one-cut containment test, not
only beyond its uniform staircase corollaries.

\paragraph{The complete support fits.}
For this partition the full nested degree-$m$ support says
that the middle forty-nine variables and the last $r-51$
variables together have degree at most $26$, while the
last variables alone have degree at most two. The exponent
of $x_2$ is limited only by the total degree $m$, because
$\lambda_2=2n>m$.
The exact number of retained indices is
\[
 G=\sum_{t=0}^2\binom{r-52+t}{t}
       \sum_{j=0}^{26-t}\binom{48+j}{j}(m-j-t+1).
                                                               \tag{O-018.5}
\]
Here $t$ is the last-variable degree, $j$ the middle degree,
and the final factor counts powers of $x_2$.
In particular,
\[
 G\leq(m+1)\binom{75}{26}\binom{r-49}{2}
          \leq\binom{75}{26}m^5<2^{75}m^5<m^{25}=n.
\]
The one-tree construction in O-018.a therefore proves
$q\leq K_n$ for this partition as well.

\paragraph{At least four actual permanent copies.}
Take the four independent highest-weight classes on the
permanent closure of degree six and weight $(6n-6,6)$,
supplied by M-012.2. Multiply all four by the common factor
\[
                   W_2^{n-16}W_{51}^{12}W_r
                                       F(e_1)^{n-r-586}.        \tag{O-018.6}
\]
All exponents are nonnegative. Its degree is
$2n-32+612+r+n-r-586=3n-6=d-6$.
After adding the seed weight, its second row is
$6+2(n-16)+24+2=2n$; rows three through fifty-one
are $24+2=26$; the remaining nonfirst rows are two.
The first row is determined by the total degree, giving
exactly (O-018.3).

O-011 supplies a nondegenerate permanent Hessian at a point
where the value of the form and its radial Hessian entry
are nonzero. Extend that radial direction to an orthogonal
basis. Every leading Hessian minor, including $W_2,W_{51},W_r$,
is then nonzero. Thus (O-018.6) is a nonzero class on the
permanent closure. Its reduced coordinate ring is a domain,
so multiplication by this one class preserves the independence
of all four seed classes. This proves the lower bound in
(O-018.4); it does not require all seed classes to be nonzero
at the particular Hessian point.

\paragraph{Every single cut is too large as a universal family.}
For $2\leq k\leq r$, let $\mathcal U_k$ be the whole
space of degree-$m$ forms whose degree in variables
$k,\ldots,r$ is at most $\min(m,\lambda_k)$, with no
other later-row restriction. Also allow $k=r+1$, with an
empty last group and no restriction. Its padded copy has
the same dimension. We show
\[
                         \dim\mathcal U_k>m^{52}
                                    >(r-2)n^2+2         \tag{O-018.7}
\]
for every cut. By O-010, none of these whole padded families
can be contained in $\Omega_n$, even allowing limits.

If $k=2$, all degrees through $m$ are allowed in the last
group. If $k>r/2$, at least $r/2-2$ early variables are
unrestricted. In either case squarefree monomials of degree
twenty-seven in an available group, with the remaining degree
on $L$, give
\[
 \dim\mathcal U_k\geq\binom{r/2-2}{27}
       \geq\frac{(r/4)^{27}}{27!}
       =\frac{m^{54}}{2^{54}27!}>m^{52}.                \tag{O-018.8}
\]
The factorial inequality follows from
$2^{54}27!<2^{54+135}<2^{400}=m^2$.

For $3\leq k\leq51$, the last-variable cap is twenty-six
and at least one early variable is available. Choose
twenty-six distinct last variables and any power through
$m-26$ of one early variable. This gives
\[
 \dim\mathcal U_k\geq(m-25)\binom{r-50}{26}
       \geq\frac{m^{53}}{2^{27}26!}>m^{52},            \tag{O-018.9}
\]
using $m-25\geq m/2$, $r-75\geq r/2$, and
$2^{27}26!<2^{157}<m$.

For $52\leq k\leq r/2$, use fifty early variables and
last-variable degree exactly two. The dimension is at least
\[
 \binom{m+48}{50}\binom{r/2+1}{2}
       \geq\frac{m^{54}}{2^{53}50!}>m^{52},            \tag{O-018.10}
\]
because $2^{53}50!<2^{353}<m^2$.
These ranges exhaust all cuts, including the empty-tail
case in (O-018.8). Finally,
$(r-2)n^2+2=(m^2-2)m^{50}+2<m^{52}$.
This proves (O-018.7).

\paragraph{Other explicit gates and scope.}
The BIP threshold holds with equality. All F-017 staircase
inequalities hold: the rows through fifty-one in question
are twenty-six and row fifty-two is two. The full-multiplicity
stabilization criteria requiring $n\geq\lambda_2$ fail
since $\lambda_2=2n$. The universal quadratic and grouped
Chow degree criteria fail: their exact thresholds at $L=dm$
are at least $\binom d2>n$.
Grenet's full containment criterion fails because
$m=2^{200}>5001$ and hence $2^m-1>2^{5000}=n$.
The older coordinate-box test fails as well: some part in
every split has at least $r/2$ variables, so its count is
at least $\binom{r/2}{27}>n$.
The coarse padding inequalities hold since
$|\bar\lambda|=2n+2r+1172\leq md$ and $r=m^2$.

The finite certificate checks the exact count (O-018.5),
the three uniform lower bounds covering every cut, the
joint-row constant in (O-018.2), and every stated numerical
comparison. No enumeration of the enormous number of cuts
is needed. This is an excluded shape with actual multiplicity
at least four, and a strict reduction of the earlier explicit
conditions. It is not an actual $q>K_n$ obstruction.

\paragraph{Geometric dictionary.}
The entire nested profile selects a downward-closed finite
torus-weight coefficient subscheme. Passing from individual
row caps to this profile is a calculation with exterior
matrix coefficients and their Cauchy span, as in F-016;
the exact background locators are Macdonald I(4.3), p.~63,
and I Appendix~A(5.3)--(5.6), pp.~155--156, and SGA~3 I,
Proposition~4.7.3, pp.~35--36. The one-tree substitution is
an integral morphism to the linear-matrix scheme. The actual
comparison is obtained only after retaining every
highest-weight ideal copy in characteristic zero. The
dimension obstruction uses the image closure, while the
four permanent copies use multiplication in its integral
coordinate ring; these are distinct arguments on the actual
schemes, not deductions from a common virtual $K_0$ class.

\section{The dimension boundary for universal jet dominance}

\subsection{O-019: the next full jet order cannot be dominant at the same power}
Let $n\geq2$, $3\leq r\leq n^2$, and $0\leq b\leq n$.
Write a degree-$n$ form in variables $x_1,z_1,\ldots,z_{r-1}$
as
\[
 F=\sum_{|\alpha|\leq n}
                  c_\alpha x_1^{n-|\alpha|}z^\alpha.
\]
Its jet of order $b$ at the first coordinate direction means
the coefficient vector $(c_\alpha)_{|\alpha|\leq b}$,
equivalently the terms of total $z$-degree at most $b$ in
$F(1,z)$. The whole jet space has dimension
\[
                    D_b^{\mathrm{jet}}(r)=\binom{r-1+b}{b}.    \tag{O-019.1}
\]
If the jets of $r$-variable size-$n$ linear determinants
are Zariski dense in this entire space, necessarily
\[
                   D_b^{\mathrm{jet}}(r)\leq(r-2)n^2+2.       \tag{O-019.2}
\]
The same necessary condition holds when one first takes
the determinant orbit closure and then projects to jets.

\begin{proof}
The jet operation is a linear projection on the affine
coefficient space. Its restriction to the closure of the
$r$-variable determinant-pencil image has image closure of
dimension at most the dimension of that source. By O-010,
the latter is at most $(r-2)n^2+2$.
The identification of this pencil closure with
$\Omega_n\cap\operatorname{Sym}^n((\mathbb C^r)^*)$
is recalled in O-017 and already includes border limits.
A dense jet image would have dimension $D_b^{\mathrm{jet}}(r)$, proving
(O-019.2). The stars-and-bars count of monomials in $r-1$
variables through degree $b$ gives (O-019.1).
\end{proof}

Now fix an integer $c\geq1$, take $r=m^2$, and let
$n=m^c$. For a fixed $b$ the two sides have orders
\[
                   D_b^{\mathrm{jet}}(m^2)=\Theta_b(m^{2b}),
                \qquad (m^2-2)m^{2c}+2=\Theta(m^{2c+2}).
\]
Thus universal jet dominance is impossible for all sufficiently
large $m$ when $b\geq c+2$. In particular, the order
$b=c+1$ reached in M-013 for $c\geq6$ is the largest
fixed integer order compatible with dimension at this power
threshold. Raising that whole jet order by one cannot be
achieved by a different matrix construction, even one that
uses limits, without increasing the power range.

For a finite BIP-scale check, take $m=10^{15}$,
$r=m^2$, $n=m^{25}$, and $b=27$. Exact arithmetic gives
\[
                \binom{m^2+26}{27}>(m^2-2)m^{50}+2.
\]
The accompanying certificate records the two integers and
their difference. This is a dimension check of the complete
jet target, not a computation of a Kronecker coefficient or
of an actual multiplicity.

\paragraph{Why this does not conflict with the positive jet theorem.}
M-013 gives full determinant multiplicity for the indicated
highest weights by proving dominance only after discarding
the higher coefficient terms. Those higher coefficients are
free in the matrices producing a prescribed low-order jet.
Consequently jet dominance does not assert an exact determinant
for each truncated padded form, and it does not imply
containment of the whole truncated coefficient linear space
in the determinant orbit closure. O-016's exact-representation
boundary can therefore coexist with order-twenty-six jet
dominance in the BIP range.

The new dimension boundary limits dominance onto the
\emph{entire} next jet space. It leaves open projections
onto smaller nested supports, interpolation on a particular
permanent coefficient image, and direct comparisons of the
two evaluation ranks. In particular it is not a proof of
an actual multiplicity obstruction beyond the stated order.

\paragraph{Geometric dictionary.}
The finite jet coordinates form a quotient of the torus-graded
coefficient module, under SGA~3, I, Proposition~4.7.3,
pp.~35--36. The jet map is a morphism of affine coefficient
schemes; dominance is the injectivity of the associated
map on coordinate algebras. Its dimension cannot exceed
that of the determinant image closure. This argument uses
neither exact lifting of all jet points nor normality of
that image. The highest-weight comparison from such a
dominant projection is the separate characteristic-zero
all-copy argument in M-013.

\section{Full multiplicity from transverse jets along a binary plane}

\subsection{O-020: arbitrary second rows with third row at most two}
Let $r\geq3$ and $n\geq r-1$. For every outer degree $d$
and every partition $\lambda\vdash nd$ with
$\ell(\lambda)\leq r$ and $\lambda_3\leq2$,
\[
                         K_n(\lambda)=a_\lambda(d[n]).
                                                               \tag{O-020.1}
\]
Missing third parts are zero. There is no bound on
$\lambda_2$ and no upper bound on $d$.
The geometric input is dominance of the determinant projection
onto the entire coefficient space of degree-$n$ forms whose
total degree outside a fixed binary plane is at most two.
Higher transverse coefficients remain unrestricted.

Write the variables as $x,y,z_1,\ldots,z_R$, where $R=r-2$.
The target coefficient space is
\[
 \mathcal T_{n,2}=
 \operatorname{Sym}^n\mathbb C^2
 \oplus\left(\operatorname{Sym}^{n-1}\mathbb C^2
                               \otimes\mathbb C^R\right)
 \oplus\left(\operatorname{Sym}^{n-2}\mathbb C^2
                               \otimes\operatorname{Sym}^2\mathbb C^R\right).
                                                               \tag{O-020.2}
\]
Its dimension is $(n+1)+Rn+\binom{R+1}{2}(n-1)$.
The map from all $n$-by-$n$ linear pencils in these variables
to $\mathcal T_{n,2}$, obtained by deleting terms of
total $z$-degree at least three from the determinant,
has a surjective differential at an explicit integral pencil.

\paragraph{A difference-labeled upper-triangular pencil.}
Since $n\geq R+1$, put
\[
 \zeta_i=z_i\quad(1\leq i\leq R),\qquad
 \zeta_i=0\quad(R+1\leq i\leq n),\qquad
 u_{ij}=\zeta_i-\zeta_j\quad(i<j).
\]
Let $U=(u_{ij})$ be strictly upper triangular, choose distinct integers
$a_1,\ldots,a_n$, and define
\[
 \ell_i=y-a_ix,\qquad \Delta=\prod_{i=1}^n\ell_i,
             A_0=\operatorname{diag}(\ell_1,\ldots,\ell_n)-U.
                                                               \tag{O-020.3}
\]
Then $\det A_0=\Delta$.

\paragraph{Binary coefficients and residue coordinates.}
The $n$ polynomials $\Delta/\ell_i$ form a basis of the
binary forms of degree $n-1$. After setting $x=1$, their
evaluation matrix at $y=a_j$ is diagonal with nonzero
diagonal entries.
Send a polynomial $p(y)$ of degree
at most $n-2$ to the vector
\[
                \left(\frac{p(a_i)}{\Delta'(a_i)}\right)_{i=1}^n.
\]
Lagrange interpolation shows that its coordinates sum to
zero; the map is injective and both spaces have dimension
$n-1$, so it identifies this polynomial space with the
zero-sum hyperplane $H_n\subseteq\mathbb C^n$. The image of
$\Delta/((y-a_i)(y-a_j))$ is
$(e_i-e_j)/(a_i-a_j)$. Thus, after rescaling each independent
edge parameter by a nonzero scalar, the binary coefficients
$\Delta/(\ell_i\ell_j)$ can be treated as the vectors
$e_i-e_j$ in $H_n$.

\paragraph{A quadratic incidence isomorphism.}
First take $n=R+1$, and let $Z_1$ denote the linear forms
in $z_1,\ldots,z_R$. Consider
\[
 \Phi:\bigoplus_{1\leq i<j\leq R+1} Z_1
       \longrightarrow H_{R+1}\otimes\operatorname{Sym}^2Z_1,
 \qquad
 (f_{ij})\longmapsto
   \sum_{i<j}(e_i-e_j)(\zeta_i-\zeta_j)f_{ij}.
                                                               \tag{O-020.4}
\]
Its source and target both have dimension
$R\binom{R+1}{2}$. The map is an isomorphism.

To prove injectivity, write $f_{ji}=f_{ij}$ and
$f_{ij}=\sum_{k=1}^R c_{ij,k}z_k$. A kernel element satisfies
\[
          \sum_{j\neq i}(\zeta_i-\zeta_j)f_{ij}=0
                                  \quad(1\leq i\leq R+1).
\]
For distinct $i,j\leq R$, the coefficient of $z_j^2$
in the $i$-th equation gives $c_{ij,j}=0$.
The $j$-th equation likewise gives $c_{ij,i}=0$.
For three distinct $i,j,k\leq R$, the coefficient of
$z_jz_k$ in the $i$-th equation gives
$c_{ij,k}+c_{ik,j}=0$. Applying the same relation at
$j$ and $k$, and using $c_{ji,k}=c_{ij,k}$, shows that
$c_{ij,k}=-c_{ij,k}$. Hence it is zero in characteristic
zero. This also covers $R=1,2$, where there are no such
remaining triple coefficients. Every $f_{ij}$ with
$i,j\leq R$ therefore vanishes. The $i$-th equation now
reduces to $z_i f_{i,R+1}=0$, which forces the remaining
linear forms to vanish. This proves injectivity and the
dimension count proves surjectivity.

For $n>R+1$ add the vertices one at a time. The previous
image is $H_{n-1}\otimes\operatorname{Sym}^2Z_1$, embedded
with last coordinate zero. For a new vertex $v$, the
edges $\{i,v\}$, $1\leq i\leq R$, have labels $z_i$.
Their new residue coordinates are $-z_i f_{iv}$, which
span every quadratic polynomial, since every quadratic
monomial is a product of a variable with a linear form.
Choose these parameters to obtain any prescribed new
coordinate. The remaining error has new coordinate zero
and total coordinate sum zero, so it belongs to the old
image and can be removed by the old edge parameters.
This proves surjectivity onto
$H_n\otimes\operatorname{Sym}^2Z_1$ for every $n\geq R+1$.
It checks every additional residue coordinate, rather than
merely comparing the numbers of parameters.

\begin{proof}[Proof of the surjective differential]
The cofactor derivative formula is
\[
             d(\det)_{A_0}(H E_{ij})
                         =H\,\Delta\,(A_0^{-1})_{ji}.
                                                               \tag{O-020.5}
\]
The expression is polynomial although it is convenient to
compute the inverse over $\mathbb C(x,y,z)$.
Since $U$ is strictly upper triangular, an inverse entry
is a finite sum over increasing paths. In particular,
\[
 (A_0^{-1})_{ii}=\ell_i^{-1},\qquad
 (A_0^{-1})_{ij}
       =\frac{u_{ij}}{\ell_i\ell_j}
                    +\text{terms of total $z$-degree at least two}
                       \quad(i<j).                    \tag{O-020.6}
\]
Every extra path edge contributes another $z$ factor.
Multiplication by $\Delta$ clears the denominators,
because a strictly increasing path visits each vertex once.

Vary diagonal entries first. Taking $H=x$ or $H=y$
in (O-020.5) spans every binary degree-$n$ form:
the polynomials $\Delta/\ell_i$ span degree $n-1$,
and multiplication by $x,y$ spans the next degree.
Taking $H=z_a$ on diagonal entries spans independently
every term linear in the $z$ variables with an arbitrary
binary coefficient of degree $n-1$.

For any edge $\{i,j\}$, $i<j$, vary the opposite,
lower-triangular entry by an arbitrary $H\in Z_1$.
Equations (O-020.5)--(O-020.6), after deleting transverse
degree at least three, give exactly
\[
                  (\zeta_i-\zeta_j)H
                          \frac{\Delta}{\ell_i\ell_j}.
                                                               \tag{O-020.7}
\]
The terms from longer paths have degree at least three
after multiplication by $H$ and disappear from the
target projection. The residue identification and the
surjectivity of (O-020.4), including its added-vertex
extension, show that these directions span every quadratic
transverse polynomial with arbitrary binary coefficients
of degree $n-2$.

These directions therefore span every summand of
(O-020.2). Diagonal directions control transverse orders
zero and one, while the lower-triangular directions just
used have no contributions in those orders. All matrix-entry coefficients are
independent parameters. Thus the differential is surjective,
with no requirement to fix higher transverse terms.
The algebraic lowest-homogeneous-term argument in M-013
now proves dominance of this coefficient projection.
\end{proof}

\begin{proof}[Proof of the multiplicity equality]
Let $w$ be any degree-$d$ highest-weight vector of weight
$\lambda$ in the determinant ideal. It depends only on
the first $r$ coordinates. By the $j=3$ support bound
of F-016, every coefficient generator used by $w$ has
total degree at most $\lambda_3\leq2$ in variables
$3,\ldots,r$. Hence $w$ is a polynomial in the
coordinates of $\mathcal T_{n,2}$.
Every linear pencil determinant belongs to $\Omega_n$,
so this polynomial vanishes after the dominant projection
just proved. Its pullback is injective, so $w=0$.
This holds for every highest-weight copy, and therefore
the ideal multiplicity space vanishes entirely. The
ambient quotient then gives (O-020.1).
\end{proof}

For permanent shapes $r\leq m^2$, the condition
$n\geq m^2-1$ is enough for this full equality whenever
$\lambda_3\leq2$, with the one- and two-row cases already
covered by O-006. This conclusion is stronger than a
comparison with a particular permanent multiplicity: every
ambient copy survives, even with an arbitrarily large second
row. It is a different sufficient condition from M-013's
bounded second-row jet theorem.

\paragraph{An independently checkable tree construction.}
The earlier proof with the weaker sufficient bound
$n\geq2R=2r-4$ is retained in
\texttt{binary\_transverse\_jets\_tree.tex}.
It colors $R$ edge-disjoint spanning trees by the individual
variables $z_a$; each tree's edge cofactors form a binary
degree-$n-2$ basis. The paths
$j,(j-i,j+i)_{1\leq i\leq R-1},j-R$ on $\mathbb Z/(2R)$
give an explicit packing, including $R=1$.
Its independent cofactor calculations pass eight exact
Jacobian rank tests, including the smallest case $R=1,n=2$.
The difference-label argument above improves the threshold
by using the quadratic incidence isomorphism instead of
disjoint color classes.

\paragraph{The exact integral incidence cokernel.}
The particular map $\Phi$ in (O-020.4), with $n=R+1$,
is defined between free abelian groups of the same rank.
Its integral cokernel is
\[
          \operatorname{coker}(\Phi_{\mathbb Z})
                      \cong(\mathbb Z/2\mathbb Z)^{\binom R3}.
                                                               \tag{O-020.8}
\]
Here the target uses the ordinary monomial basis of
$\operatorname{Sym}^2\mathbb Z^R$. Delete the last residue
coordinate to identify $H_{R+1,\mathbb Z}$ with $\mathbb Z^R$.
The $R^2$ coefficients of the edges from $i$ to the zero
vertex occur with coefficient one in the rows given by
vertex $i$ and a monomial $z_i z_k$. They occur in no
other rows after the last residue coordinate is deleted.
Unimodular elimination removes these columns and rows.

Next, the two endpoint coefficients of each finite edge
$\{i,j\}$ occur with coefficient minus one in the remaining
rows $(i,z_j^2)$ and $(j,z_i^2)$, respectively.
They are unit pivots and can also be eliminated.
The remaining coordinates split by unordered triples
$\{i,j,k\}$. Each triple has the three variables
$c_{ij,k},c_{ik,j},c_{jk,i}$ and the three relations
\[
             a+b,\qquad a+c,\qquad b+c,
\]
up to signs. This $3$-by-$3$ block has Smith form
$\operatorname{diag}(1,1,2)$: a unit entry and a unit
$2$-by-$2$ minor give its first two invariant factors,
and its determinant has absolute value two. This proves
(O-020.8), including $R=1,2$, where there are no triples.
Thus this selected quadratic incidence submap is an
isomorphism after inverting two, and its rank drops by
$\binom R3$ after reduction modulo two.
No failure of the full determinant jet map in characteristic
two is inferred; other coefficient directions are not
included in $\Phi$.

\paragraph{Geometric dictionary.}
The target is the degree-$n$ part of the quotient by
$(z_1,\ldots,z_R)^3$, the second infinitesimal neighborhood
transverse to the binary coordinate plane. It is a finite
locally free coefficient module. The difference-labeled graph, diagonal
linear forms, and cofactor derivatives define integral
morphisms. Distinctness of the $a_i$ and the nonzero
Vandermonde factors and the precisely computed two-torsion
in (O-020.8) explain the passage to
$\mathbb Q$ or $\mathbb C$; no equal-rank assertion in
every residue characteristic is inferred.
The suffix support is a torus/comodule calculation under
SGA~3 I, Propositions~4.7.2--4.7.3, pp.~34--36,
and the Cauchy and Schur functors of Macdonald I(4.3),
p.~63, and Appendix~A(5.3)--(5.6), pp.~155--156.
Dominance injects the actual coordinate algebra of the
detected coefficients, which makes the whole highest-weight
ideal component vanish. Neither exact lifting of each
truncated form nor containment of that entire linear space
in the determinant closure is asserted.

\section{Full permanent multiplicities for a thin third row}

\subsection{O-021: a transverse permanent jet and padding theorem}
Let $r\geq3$, $n\geq m$, and let $\lambda\vdash nd$ satisfy
\[
       \ell(\lambda)\leq r,\qquad \lambda_3\leq2,\qquad
                    m\geq\max\{r,\lambda_2+1\}.
\]
Then, with the same internal padding as the supplied IP and
BIP papers,
\[
                    q^m_\lambda(d[n])
                         =K_n(\lambda)=a_\lambda(d[n]).
                                                               \tag{O-021.1}
\]
There is no upper bound on the outer degree $d$.
The statement determines all copies of these types on the
permanent side, rather than only their occurrence.

\paragraph{The unpadded permanent has the same transverse differential.}
For any $s\geq r-1$, use the variables
$x,y,z_1,\ldots,z_R$, $R=r-2$, and the difference-labeled
matrix of O-020, but change the signs on its upper-triangular
entries:
\[
        A_0=\operatorname{diag}(\ell_1,\ldots,\ell_s)+U,
        \qquad \ell_i=y-a_i x,\qquad
        U_{ij}=\zeta_i-\zeta_j\quad(i<j).
\]
Here $\zeta_i=z_i$ for $i\leq R$ and zero otherwise,
and the $a_i$ are distinct. The underlying directed graph
is acyclic. Consequently
$\operatorname{per}A_0=\Delta=\prod_i\ell_i$.

The permanent minor obtained by deleting row $j$ and
column $i$, with $i<j$, is the sum over directed paths
from $i$ to $j$ of the path's edge product multiplied
by the diagonal entries at its unused vertices.
Indeed every term of that minor consists of this path
and loops; an additional nontrivial cycle is impossible
in an acyclic graph. Conversely each path and its unused
loops gives exactly one minor term. In particular the
minor equals
\[
        (\zeta_i-\zeta_j)\frac{\Delta}{\ell_i\ell_j}
               +\text{terms of transverse degree at least two}.
                                                               \tag{O-021.2}
\]
The diagonal minors are $\Delta/\ell_i$.
This is an identity with weighted paths; negative coefficients
inside the difference labels cause no difficulty.

Vary diagonal entries by $x,y,z_a$, and lower-triangular
entries by linear forms in $z$. After deleting transverse
degree at least three, their permanent differentials are
exactly the selected determinant differentials in O-020.
The quadratic incidence isomorphism and the added-vertex
argument therefore apply unchanged. The permanent map
from all $s$-by-$s$ linear pencils dominates the full space
\[
 \operatorname{Sym}^s\mathbb C^2
 \ \oplus\ \bigl(\operatorname{Sym}^{s-1}\mathbb C^2
                                      \otimes\mathbb C^R\bigr)
 \ \oplus\ \bigl(\operatorname{Sym}^{s-2}\mathbb C^2
                               \otimes\operatorname{Sym}^2\mathbb C^R\bigr).
                                                               \tag{O-021.3}
\]
Its higher transverse coefficients are unrestricted.

\begin{proof}[Proof of the full multiplicity equality]
If $a_\lambda=0$, both quotient multiplicities vanish,
so assume it is positive. One-row types have
$q=K=a=1$, since the highest-weight function $F(e_1)^d$
is nonzero on both closures. Thus assume $\lambda_2>0$.
Then $\ell(\lambda)\leq d$.
Set $s=m-1$ and $b=\lambda_2$. The hypotheses give
$s\geq r-1$ and $s\geq b$. By F-014, the complete
ambient highest-weight multiplicity spaces at inner degrees
$s$ and $n$ have the same dimension:
\[
 a_\mu(d[s])=a_\lambda(d[n]),\qquad
             \mu=(sd-|\bar\lambda|,\bar\lambda).
                                                               \tag{O-021.4}
\]
Here $\bar\lambda=(\lambda_2,\lambda_3,\ldots)$.
The first part of $\mu$ is at least its second part,
because $|\bar\lambda|\leq(d-1)b$ and $s\geq b$.

The padding lift $T$ of O-012.b is injective from the
entire highest-weight space of weight $\mu$ into that
of weight $\lambda$. Equation (O-021.4) makes it an
isomorphism of these multiplicity spaces.
For the actual padding map
$P(F)=x^{n-s}F$, its pullback on $T$ is the diagonal
operator from the differential pairing in O-012.b:
\[
 P^*T=\operatorname{Sym}^d(\mathcal D),\qquad
 \mathcal D(e^\alpha)
        =\frac{(\alpha_1+n-s)!}{\alpha_1!}\,e^\alpha,
                 \qquad |\alpha|=s.                   \tag{O-021.5}
\]
Every diagonal entry is nonzero in characteristic zero.
The operator also preserves every monomial support condition.
It need not preserve the highest-weight subspace; that
additional property is not used.

By F-016 at its third-row cut, every polynomial in the
highest-weight space of weight $\mu$ uses only coefficient
generators of transverse degree at most two, since
$\mu_3=\lambda_3\leq2$. Thus the same is true after
applying the invertible operator in (O-021.5).
Any nonzero linear combination of the lifted copies
therefore pulls back to a nonzero polynomial in the
coordinates of (O-021.3). The dominant permanent map
proved above cannot annihilate that polynomial.
This proves independence of every lifted copy on the
padded permanent-pencil family.

Finally that entire family belongs to the required
internally padded permanent closure. For an arbitrary
$s$-by-$s$ pencil $A$, use the $m$-by-$m$ block matrix
\[
                         \widetilde A=\operatorname{diag}(x,A).
\]
Its distinguished entry is exactly $\widetilde A_{11}=x$,
and
\[
 \widetilde A_{11}^{\,n-m}\operatorname{per}\widetilde A
                      =x^{n-s}\operatorname{per}A.
                                                               \tag{O-021.6}
\]
These are linear endomorphism substitutions in
$X_{11}^{n-m}\operatorname{per}_m$. They belong to its
orbit closure because invertible endomorphisms are dense.
The variable space fits: $r\leq m\leq m^2\leq n^2$.
The independence just proved forces the entire permanent
ideal highest-weight space of type $\lambda$ to vanish.
Hence $q^m_\lambda=a_\lambda$.
O-020 also gives $K_n=a_\lambda$, since $n\geq m\geq r$.
This proves (O-021.1), including $n=m$.
\end{proof}

\paragraph{A strict four-copy example below every exact support-cover size.}
For a downward-closed support $\Gamma$ containing
$x_2^b$ and $x_i^2$ for every $3\leq i\leq r$, with
$b\geq2$, the two-tree invariant of F-018 satisfies
\[
                            \kappa(\Gamma)\geq b+r-2.
                                                               \tag{O-021.7}
\]
To prove this, let $P,Q$ be any admissible pair.
A decomposition of the pure power $b e_2$ forces
$a e_2\in P$, $a'e_2\in Q$ with $a+a'\geq b-1$.
Downward closure supplies both full axis chains, so these
alone contribute at least $a+a'+1\geq b$ to
$|P|+|Q|-1$. For $2e_i$, at least one of the two
endpoints has a positive $i$-th coordinate: the coefficient
edge can contribute at most one such unit. Downward closure
then forces $e_i\in P\cup Q$. These $r-2$ distinct
off-axis vertices are additional to the two chains.
This proves (O-021.7).

Take the concrete values
\[
       m=9,\quad n=10,\quad d=r=6,\qquad
                    \lambda=(44,8,2,2,2,2).
                                                               \tag{O-021.8}
\]
O-021.1 applies because $m=\max\{6,8+1\}$.
The exact near-rectangle formula F-013 gives
\[
      a_\lambda(6[10])=p_6(6,8)-p_5(6,8)=11-7=4.
\]
Thus
\[
                        q^9_\lambda(6[10])
                                =K_{10}(\lambda)=4.   \tag{O-021.9}
\]
Its nested support is
$\Gamma=\{\gamma:|\gamma|\leq8,\
                  \gamma_3+\cdots+\gamma_6\leq2\}$.
It contains the powers required by (O-021.7), so
$\kappa(\Gamma)\geq12>n>m$. Since the total-degree
cap is already eight, this is the same support at base
degrees $8,9,10$. Every exact two-tree support-cover
criterion of F-018 therefore misses this example, while
the transverse dominant projection and the all-copy lift
determine its actual multiplicities. This is a limitation
of that exact-cover architecture, not a lower bound for
arbitrary determinant or permanent representations.

\paragraph{Geometric and integral scope.}
The weighted path minors and the block substitution
(O-021.6) are integral polynomial identities. The
transverse differential uses the nonzero Vandermonde
factors and the selected incidence map whose exact
two-torsion is computed in O-020.8. The padding pullback
also uses invertible factorial scalars in characteristic
zero. No equality of integral Schur decompositions after
arbitrary base change is asserted.
The torus support is the comodule grading of SGA~3 I,
Propositions~4.7.2--4.7.3, pp.~34--36; F-014 and F-016
use the Cauchy and exterior functors of Macdonald I(4.3),
p.~63, Appendix~A(5.3)--(5.6), pp.~155--156, and
SGA~6 VI, Theorem~3.3, pp.~392--393.
The argument retains the injective coordinate pullback
through the dominant permanent projection and the actual
padding map. It supplies multiplicity information beyond
the virtual lambda class, with no saturation inference.

\section{A positive BIP example detected beyond exact support covers}

\subsection{O-022: the prefix jet criterion makes a strict further reduction}
Set
\[
 m=2^{200},\qquad r=m^2,\qquad n=m^{25}=2^{5000},
 \qquad d=3n,
\]
and put
\[
           \lambda=(3n^2-26r+20,\ 32,\ 26^{r-2}).
                                                               \tag{O-022.1}
\]
This is a partition of $nd$ with $r$ parts. It satisfies
\[
                  4\leq q^m_\lambda(d[n])
                                \leq K_n(\lambda)=a_\lambda(d[n]).
                                                               \tag{O-022.2}
\]
The exact values of these multiplicities are not asserted.
The full multiplicity conclusion follows from M-017.
Every exact two-tree support cover of F-018 has size greater
than $n$ for this shape. Moreover every entire one-cut
universal family is too large, even in border dimension,
and the entire point-jet target is also too large.
Thus (O-022.2) is a strict reduction beyond those proved
criteria, with actual permanent multiplicity at least four.

\paragraph{The prefix extension fits.}
Use M-017 with $b=26$, $\ell=12$, $t=6$, and $k=3$.
Here $p=1$, $\lceil\sqrt{r-1}\rceil=m$, and the size is
\[
                       N=8m\binom{m^2+11}{12}.
                                                               \tag{O-022.3}
\]
The exact integer constant in M-017 is
\[
 C_*=\left\lfloor\frac{190^{24}}{\binom{36111}{12}}\right\rfloor
                  =478126755.
                                                               \tag{O-022.4}
\]
The saved certificate gives the quotient and positive
remainder and verifies the defining two integer inequalities.
The normalized binomial ratio is nonincreasing in $m$.
Since $8<C_*$ and $m\geq190$, (O-022.3) is strictly
less than $m^{25}=n$. The target has total nonfirst
degree at most $32$ and degree at most $26$ in variables
$3,\ldots,r$. These are exactly the relevant two nested
bounds of $\lambda$. M-017 therefore gives $K_n=a_\lambda$.

\paragraph{Four permanent copies survive.}
M-012 supplies four independent permanent-closure
highest-weight classes of degree six and weight
$(6n-6,6)$. Write $c_0(F)=F(e_1)$ and let $W_r$
be the leading Hessian determinant from O-011.
Multiply each class by the common highest-weight function
\[
                          W_r^{13}c_0^{\,d-6-13r}.
                                                               \tag{O-022.5}
\]
Its exponent is nonnegative. O-011 gives a permanent
point and basis where $W_r$ and $c_0$ are simultaneously
nonzero, since $r=m^2$. The coordinate algebra of the
irreducible permanent orbit closure is a domain, so
multiplication by (O-022.5) preserves all four independent
classes. Their total degree is $d$; their first weight is
$nd-26r+20$, their second is $6+26=32$, and every
remaining weight is $26$. Hence $q^m_\lambda\geq4$.
The ambient quotient bound and the preceding equality
give (O-022.2).

\paragraph{Every exact support cover is too large.}
Let $\Gamma$ denote the full nested support at base degree
$32$ or at base degree $m$. Both supports are the same:
\[
       \Gamma=\{\gamma:|\gamma|\leq32,\
                            \gamma_3+\cdots+\gamma_r\leq26\}.
                                                               \tag{O-022.6}
\]
It contains every homogeneous degree-$26$ monomial in
the $R'=r-2$ last variables. If a support cover had
size $S\leq n$, its integral determinant identity would
represent every such form $f$ as
$x_1^{S-26}f$. Here $S\geq26$ by F-018.
Specializing $x_1=1$ and the unused early variable $x_2=0$
would give an affine linear determinant of size at most
$n$ for every homogeneous degree-$26$ form in $R'$ variables.

This contradicts the exact generic lower bound O-016.
Indeed its middle-layer ratio is
\[
 \frac{D_{26}(R')}{2D_{13}(R')}
     =\frac{13!}{2\cdot26!}\prod_{j=13}^{25}(R'+j)
     \geq\frac{13!\,m^{26}}{2^{14}\,26!}>m^{25}=n.
                                                               \tag{O-022.7}
\]
We used $R'=m^2-2\geq m^2/2$ and
$2^{14}26!/13!<2^{79}<m$. Thus a generic form has
exact determinantal complexity greater than $n$, and
$\kappa(\Gamma)>n$. This argument concerns exact
representations. It does not infer a border-complexity
lower bound for this restricted family.

\paragraph{Every one-cut enlargement fails even by dimension.}
The determinant-pencil closure in $r$ variables has
dimension at most
\[
                      (r-2)n^2+2<m^{52}.
                                                               \tag{O-022.8}
\]
For a cut $k$, the entire universal degree-$m$ family
allows unrestricted early degree in variables
$2,\ldots,k-1$ and suffix degree at most $\lambda_k$,
with total degree at most $m$.

At $k=2$, it contains all squarefree degree-$27$
monomials in the last variables, because $\lambda_2=32$.
At every cut $k>r/2$, including $k=r+1$, there are at
least $r/2-1$ unrestricted early variables. Choosing
$27$ of them gives in both cases the common lower bound
\[
     \binom{r/2-1}{27}
          \geq\frac{m^{54}}{2^{54}27!}>m^{52}.
                                                               \tag{O-022.9}
\]
The last inequality follows from
$2^{54}27!<2^{189}<m^2$.

For $3\leq k\leq r/2$, there are at least $r/2$
suffix variables, and their allowed total degree is
$26$. Choose $26$ distinct suffix variables and an
arbitrary exponent from $0$ through $m-26$ on $x_2$.
The family therefore has dimension at least
\[
 (m-25)\binom{r/2}{26}
           \geq\frac{m^{53}}{2^{53}26!}>m^{52},
                                                               \tag{O-022.10}
\]
using $m-25\geq m/2$ and
$2^{53}26!<2^{183}<m$.
These three ranges exhaust all cuts. Multiplication by
the padding power is injective, so the same dimensions
hold for the padded linear spaces. Equations
(O-022.8)--(O-022.10) exclude containment in the
determinant closure, including border containment.
The full point-jet target through degree $32$ contains
the monomials counted in (O-022.9), so universal dominance
onto that entire target is likewise impossible by O-019.
The restricted projection of M-017 avoids this dimension
obstruction by retaining both nested bounds.

\paragraph{The preceding explicit candidate conditions are met.}
The shape satisfies $\ell(\lambda)=m^2$,
$|\bar\lambda|=26m^2-20\leq md$, and
$\ell(\lambda)\leq d$. It meets the complete BIP
staircase and the joint A-025 condition, since
$\lambda_2=32$, all subsequent parts are $26$, and
$\lambda_{52}=26$. It is outside O-020's
third-row-at-most-two range.
The uniform exact Chow degree threshold $P_d(dm)$ of
A-023 is at least $\binom d2>n$ here.
Its stronger shape-specific threshold also fails:
$P_2(L)=1$ and every recurrence step increases by at
least one for $L\geq1$, so
$P_d(32)\geq d-1=3n-1>n$.
The whole-permanent Grenet threshold also fails:
$2^m-1>2^{5000}=n$.

For completeness, every coordinate-box split of M-009
has a group containing at least $r/2$ variables.
All individual caps are at least $26$, so it contains
all squarefree monomials of degree $m$. Its size is at least
\[
                 \binom{r/2}{m}\geq(m/2)^m>m^{25}.
                                                               \tag{O-022.11}
\]
The binomial inequality follows termwise from
$(r/2-i)/(m-i)\geq(r/2)/m$.
The final comparison follows from $m=2^{200}$.
All exact support-cover variants are already excluded
more generally by (O-022.7).
These checks establish strict refinement of the listed
proved conditions. They do not exhibit an actual
multiplicity obstruction: (O-022.2) proves the opposite
comparison for this particular shape.

\paragraph{Geometric dictionary and exact versus border scope.}
The target (O-022.6) is a finite torus-weight coefficient
module under SGA~3 I, Proposition~4.7.3, pp.~35--36.
M-017 gives a universally injective pullback on its
polynomial coordinate algebra, and F-016 identifies every
relevant Schur-copy coefficient using the exterior functors
of SGA~6 VI, Theorem~3.3, pp.~392--393, and Macdonald
I(4.3), p.~63, Appendix~A(5.3)--(5.6), pp.~155--156.
The four copies in (O-022.5) are multiplied in the actual
integral coordinate algebra of an irreducible orbit closure,
not only in its virtual representation class.
Equation (O-022.7) uses an exact-complexity bound and is
kept separate from the border-valid dimension statements
(O-022.8)--(O-022.10). The latter use the actual reduced
image closure and its projection, as in O-010.
Higher coefficients remain free in the dominant prefix
projection; no exact lift of every truncated form is asserted.

\section{A further positive example requiring a nonsimplex extension}

\subsection{O-023: the order-ideal extension improves every prefix-simplex test}
Keep
$m=2^{200}$, $r=m^2$, $n=m^{25}$, and $d=3n$,
and now define
\[
 \lambda=(nd-26r-1000000+50,\ 1000000,\ 28,\ 26^{r-3}).
                                                               \tag{O-023.1}
\]
Then
\[
                 4\leq q^m_\lambda(d[n])
                               \leq K_n(\lambda)=a_\lambda(d[n]).
                                                               \tag{O-023.2}
\]
Every choice of parameters in the prefix-simplex criterion
M-017 fails to reach this shape at size $n$.
The exact support-cover and one-cut failures of O-022
also remain. Thus M-018 makes a further strict reduction
of the proved candidate conditions.
No exact numerical value of the multiplicities in
(O-023.2) is asserted.

\paragraph{The nonsimplex extension fits.}
For the order ideal
\[
 A=\{x_2^a x_3^c:\ a+c\leq999974,\quad c\leq28\}
\]
one has
\[
       |A|=29\cdot999975-\frac{28\cdot29}{2}
                          =28998869.
                                                               \tag{O-023.3}
\]
The full nested support of (O-023.1) has total degree
at most one million, suffix degree from coordinate three
at most twenty-eight, and suffix degree from coordinate
four at most twenty-six.
If its monomial degree exceeds twenty-six, choose an
early submonomial in $x_2,x_3$ of the excess degree.
There are enough early factors because the later suffix
has degree at most twenty-six. Its degree is at most
$999974$, and its $x_3$ exponent is at most twenty-eight,
so it lies in $A$. The remaining monomial has degree
twenty-six. Smaller monomials use the zero element of $A$.
Thus the full support is contained in $A+D_{26}$.

Use M-018 with $b=26$ and
$T=m\binom{m^2+11}{12}$.
The size is $28998870T<n$, because
$28998870<C_*=478126755$ and the uniform normalized
estimate of M-017 applies. This proves $K_n=a_\lambda$.

\paragraph{Four positive permanent copies.}
Multiply the four degree-six binary seed classes of M-012
by
\[
             W_r^{13}\,W_3\,W_2^{499983}\,
                         c_0^{\,d-13r-999975}.          \tag{O-023.4}
\]
The exponent is nonnegative. O-011 supplies one basis
at a permanent point where every displayed leading
Hessian minor and $c_0$ is nonzero. Therefore the common
factor is nonzero in the integral coordinate algebra of
the irreducible permanent closure, and multiplication
preserves the four independent seed classes.
Their degree is
$6+13r+3+2\cdot499983+(d-13r-999975)=d$.
Their second weight is
$6+26+2+999966=1000000$, their third weight is $28$,
and all later weights are $26$. The first weight is
the remaining part $nd-26r-999950$, as required.
This proves $q^m_\lambda\geq4$ and hence (O-023.2).

\paragraph{All prefix-simplex choices fail.}
Consider any M-017 parameter choice that could cover this
highest-weight support. Its even base order is $b=2\ell+2$,
and it must satisfy $\lambda_2\leq b+t$ and
$\lambda_k\leq b$. Taking more ambient variables than
the actual length $r$ only increases its size; it suffices
to give lower bounds using this length.

If $b\geq28$, then $\ell\geq13$, and even the unextended
core has size at least
\[
     2m\binom{r+12}{13}
           \geq\frac{2m^{27}}{13!}>m^{25}=n.             \tag{O-023.5}
\]
If $b=26$, the first possible cut is $k=4$, since
$\lambda_3=28$. Therefore $p=k-2\geq2$ and
$t\geq999974$. Its size is at least
\[
 \frac{m^{25}}{12!}
       \left(1+\binom{999976}{2}\right)>m^{25},
                                                               \tag{O-023.6}
\]
because $T\geq m^{25}/12!$ and the parenthesized integer
is larger than $12!$.
Finally, if $b\leq24$, every actual part is larger than
$b$, so the cut must be at or after $r+1$.
Then $p\geq r-1$, $t\geq13$, and $T\geq m$.
Its size is at least
\[
                    m\binom{r+12}{13}
                          \geq\frac{m^{27}}{13!}>n.    \tag{O-023.7}
\]
These cases exhaust all allowed even base orders.
The strict inequalities use $m^2>13!$, and are verified
by the exact certificate. Thus no M-017 parameter choice
satisfies its size condition.

\paragraph{The earlier failures persist.}
The nested support (O-022.6) is contained in the present
support: its total degree is at most $32$, its suffix
from the third coordinate is at most $26$, and therefore
its suffix from the fourth coordinate is at most $26$.
In particular the present support contains all degree-$26$
forms in the last $r-2$ variables. The exact-complexity
argument (O-022.7) again forces every F-018 exact support
cover to have size greater than $n$.

Each part of (O-023.1), except the first, is at least
the corresponding part of (O-022.1).
Every one-cut family and coordinate box is consequently
an enlargement of its O-022 counterpart. The border-valid
dimension failures (O-022.8)--(O-022.10) and the box
size failure (O-022.11) remain valid.
The full point-jet target is still too large, since it
contains the degree-$27$ monomials used there.
The shape meets the BIP staircase, A-025's joint gate,
and the M-017 coupled necessary conditions; the complete
M-017 size failure was proved above.
Its tail is $26r+999950\leq md$ and its length is $m^2$.
The shape-specific Chow threshold is at least
$d-1=3n-1>n$; the uniform threshold and Grenet failure
are unchanged. O-020 does not apply because $\lambda_3=28$.

These checks concern exact proved exclusions, including
ones with actual multiplicity at least four. They do not
claim that (O-023.1) is an obstruction: (O-023.2)
proves that it is not.

\paragraph{Geometric dictionary.}
The set $A$ is a finite monomial order ideal in an integral
torus-graded coefficient module, as in SGA~3 I,
Proposition~4.7.3, pp.~35--36. M-018 gives its actual
universally injective coordinate pullback. The full
highest-weight support is the Schur/exterior calculation
of F-016, with SGA~6 VI, Theorem~3.3, pp.~392--393,
and Macdonald I(4.3), p.~63, Appendix~A(5.3)--(5.6),
pp.~155--156. The domain multiplication retains four
sections of the actual orbit-closure coordinate algebra.
The exact-complexity obstruction and the separate
border dimension arguments retain the distinct scopes
spelled out in O-022. No saturation statement follows
from a virtual class alone.

\section{A canonical positive count for order-ideal jet extensions}

\subsection{O-024: subtracting the jet order from every suffix bound}
Let $\lambda$ have length at most $r\geq2$, and let
$b\geq0$ be an integer. Set
\[
           \mu_j=(\lambda_j-b)_+\quad(2\leq j\leq r),
            \qquad u_+=\max(u,0),
\]
and write $\Gamma_\lambda$ for the full nested suffix
support of M-018. Then
\[
                         \Gamma_\lambda\subseteq\Gamma_\mu+D_b.
                                                               \tag{O-024.1}
\]
Moreover
\[
        |\Gamma_\mu|
            =\#\{\nu:\nu\text{ is a partition and }
                         0\leq\nu_j\leq\mu_j\ (2\leq j\leq r)\}.
                                                               \tag{O-024.2}
\]
Thus the needed extension size has a direct positive count:
the number of subpartitions of the tail obtained by
subtracting $b$ and replacing negative entries by zero.
This is an explicit consequence of the general M-018
criterion, not a stronger existential criterion.

\begin{proof}
Take $\gamma\in\Gamma_\lambda$. Remove at most $b$
variable factors, always starting with the largest index,
and call their multiset $\eta$. If $|\gamma|\leq b$,
remove all factors. Set $\alpha=\gamma-\eta$.
For every suffix beginning at $j$, this procedure removes
exactly the smaller of $b$ and that suffix's original
size before removing any earlier variable. Consequently
\[
         \sum_{i=j}^r\alpha_i
                =\left(\sum_{i=j}^r\gamma_i-b\right)_+
                \leq(\lambda_j-b)_+=\mu_j.
                                                               \tag{O-024.3}
\]
Hence $\alpha\in\Gamma_\mu$ and $\eta\in D_b$,
proving (O-024.1).

For (O-024.2), send $\alpha\in\Gamma_\mu$ to its
suffix cumulative sums
$\nu_j=\sum_{i=j}^r\alpha_i$.
They form a weakly decreasing sequence of nonnegative
integers bounded componentwise by $\mu$.
Conversely any such subpartition gives
$\alpha_j=\nu_j-\nu_{j+1}$, with $\nu_{r+1}=0$.
These two operations are inverse, proving the positive count.
\end{proof}

For an even jet order $b=2\ell+2$, put
\[
 T=\lceil\sqrt{r-1}\rceil\binom{r-1+\ell}{\ell},
 \qquad c(\mu)=\#\{\nu:\nu\subseteq\mu\}.
\]
M-018 now gives the explicit sufficient condition
\[
 n\geq T(1+c(\mu))
       \quad\Longrightarrow\quad K_n(\lambda)=a_\lambda(d[n]).
                                                               \tag{O-024.4}
\]
If $n\geq m\geq T(1+c(\mu))$, it also gives
$q^m_\lambda=K_n=a_\lambda$ with internal padding.
These are full-copy conclusions at arbitrary outer degree.

\paragraph{An exact positive recurrence.}
The count requires no cancellation. For $0\leq t\leq\mu_r$
put $F_r(t)=1$, and put it zero otherwise.
Recursively define
\[
 F_j(t)=\mathbf1_{\{0\leq t\leq\mu_j\}}
                         \sum_{s=0}^t F_{j+1}(s)
                    \quad(j=r-1,\ldots,2).
                                                               \tag{O-024.5}
\]
Here $F_j(t)$ counts allowed suffixes with cumulative
sum $t$ at position $j$; choosing the next cumulative
sum $s\leq t$ proves the recurrence.
Finally $c(\mu)=\sum_{t=0}^{\mu_2}F_2(t)$.
This is a finite positive computation, not a claim of
polynomial complexity in a binary encoding of the row lengths.

\paragraph{A sharper explicit three-row condition in BIP range.}
Every actual obstruction at $n\geq m^{25}$ satisfies
\[
 \boxed{\lambda_2\geq159375611\quad\text{or}\quad
          \lambda_3\geq29\quad\text{or}\quad\lambda_4\geq27.}
                                                               \tag{O-024.6}
\]
Indeed suppose the three inequalities fail. With $b=26$,
the shaved tail has at most two nonzero parts, bounded
by $(159375584,2)$. If a two-part partition is $(A,C)$,
with $A\geq C$, its number of subpartitions is
\[
                  \sum_{j=0}^C(A-j+1)
                 =(C+1)(A+1)-\frac{C(C+1)}2.
                                                               \tag{O-024.7}
\]
Thus here
$c(\mu)\leq478126752$ and
$1+c(\mu)\leq478126753<C_*=478126755$.
The uniform normalized estimate of M-017 makes
$T(1+c(\mu))<m^{25}\leq n$, so (O-024.4)
gives $K_n=a_\lambda$ and excludes an actual obstruction.
The exact arithmetic also gives
$3\cdot159375611-77=478126756>C_*$;
the displayed constant is the largest obtained by this
particular two-part count with the stated other caps.
It is not optimal among all order-ideal choices.

\paragraph{The canonical order ideal need not be smallest.}
For the full degree-four simplex in two variables, with
$b=2$, the canonical shaved ideal is the degree-two
simplex and has six members. The five-member order ideal
\[
                     \{1,x,x^2,y,y^2\}
\]
also covers the degree-four simplex after adding $D_2$:
a monomial of degree at most two uses the member $1$;
one of degree three or four has some exponent at least
two, so remove $x^2$ or $y^2$.
Thus (O-024.2) is an explicit sufficient count, not
the solution of the finite optimization in M-018.

\paragraph{Grothendieck dictionary.}
The exponent sets are finite order ideals in a torus
character monoid, and their free coefficient modules
are comodules as in SGA~3 I, Proposition~4.7.3,
pp.~35--36. The positive count (O-024.2) indexes
an explicit integral monomial basis. M-018 supplies
the actual universally injective coordinate pullback;
F-016's Schur/exterior support calculation uses SGA~6
VI, Theorem~3.3, pp.~392--393, and Macdonald I(4.3),
p.~63, Appendix~A(5.3)--(5.6), pp.~155--156.
The partition bijection gives the size of a coefficient
module. It does not by itself give positivity of a new
Kronecker coefficient or saturation of its support;
the multiplicity conclusion still uses the proved
dominant coefficient morphism.

\section{A positive example beyond every one-sided jet extension}

\subsection{O-025: strictness of the two-sided order-ideal criterion}
Set
\[
 m=2^{200},\qquad r=m^2,\qquad n=m^{25},\qquad d=3n,
             \qquad B_0=1000000,
\]
and define
\[
 \lambda=(nd-26r-1999928,\ B_0+6,\ B_0,\ 26^{r-3}).
                                                               \tag{O-025.1}
\]
Then
\[
                  4\leq q^m_\lambda(d[n])
                                \leq K_n(\lambda)=a_\lambda(d[n]).
                                                               \tag{O-025.2}
\]
The two-sided construction M-019 proves the ambient
determinant equality. No choice of base jet order and
finite one-sided order ideal in M-018 can prove that
equality at size $n$ for this support. This compares
all of M-018's admissible choices, not just a particular
canonical ideal or a prefix simplex.
The exact values of the multiplicities remain unspecified.

\paragraph{The two-sided extension fits.}
Use the canonical shaving O-024 with $b=26$.
The shaved support has only two nonzero bounds:
\[
                       (B_0-20,\ B_0-26).
\]
It is contained in the sum of the two axis ideals
\[
 \mathcal A=\{1,x_2,\ldots,x_2^{B_0-20}\},\qquad
 \mathcal B=\{1,x_3,\ldots,x_3^{B_0-26}\}.
\]
Therefore $\Gamma_\lambda\subseteq\mathcal A+\mathcal B+D_{26}$.
At this perfect-square length, both
$\lceil\sqrt r\rceil$ and $\lceil\sqrt{r-1}\rceil$ equal $m$.
With
$T=m\binom{r+11}{12}$, the two-sided size is
\[
          N=(|\mathcal A|+|\mathcal B|)T
                =(2B_0-44)T=1999956T<n.                \tag{O-025.3}
\]
The final inequality follows from $1999956<C_*=478126755$
and M-017's normalized estimate. M-019 gives $K_n=a_\lambda$.
No uniform dominance relation between the two criteria
at other, nonsquare ambient lengths is needed here.

\paragraph{Four positive permanent copies.}
Multiply the four degree-six binary classes of M-012 by
\[
                      W_r^{13}\,W_3^{499987}\,
                             c_0^{\,d-13r-1499967}.     \tag{O-025.4}
\]
The last exponent is nonnegative, and the Hessian argument
O-011 makes all factors simultaneously nonzero on the
permanent closure. Its coordinate algebra is an integral
domain, so all four classes remain independent.
The total degree is
$6+13r+3\cdot499987+(d-13r-1499967)=d$.
The second weight is
$6+26+2\cdot499987=1000006$, the third is
$26+2\cdot499987=1000000$, and the later weights are $26$.
Their first weight is the remaining part in (O-025.1).
This proves $q^m_\lambda\geq4$ and hence (O-025.2).

\paragraph{Excluding every one-sided choice.}
Suppose, for a contradiction, that M-018 applies with
some even base order $b=2\ell+2$ and a nonempty finite
order ideal $A$:
\[
 \Gamma_\lambda\subseteq A+D_b,\qquad
             T_b(1+|A|)\leq n,\qquad
             T_b=m\binom{r-1+\ell}{\ell}.
                                                               \tag{O-025.5}
\]
Taking more ambient variables can only increase this
core size. Monomials using those additional variables
cannot participate in a nonnegative decomposition of
an exponent supported in the first $r$ variables, so
discarding them from $A$ retains every decomposition
needed below. Thus (O-025.5) is the weakest relevant
size bound even if a larger ambient length was chosen.

If $b\geq28$, then $\ell\geq13$ and
\[
 T_b(1+|A|)\geq2m\binom{r+12}{13}
                       \geq\frac{2m^{27}}{13!}>n.
                                                               \tag{O-025.6}
\]
This already contradicts feasibility.

If $b=26$, then $T_b\geq n/12!$, so
$|A|<12!=479001600$. There are $r-3>12!$ suffix
variables $x_4,\ldots,x_r$. At least one of their unit
monomials, say $x_j$, is missing from $A$.
Downward closure implies that no member of $A$ uses $x_j$.
For every two-variable monomial $\nu(x_2,x_3)$ of total
degree at most $B_0-26=999974$, the monomial
$x_j^{26}\nu$ lies in $\Gamma_\lambda$.
In a decomposition of its exponent as $\alpha+\eta$,
with $\alpha\in A$ and $|\eta|\leq26$,
the entire $x_j^{26}$ must belong to $\eta$.
It exhausts that degree budget, so $\eta=26e_j$ and
$\alpha=\nu$. Consequently $A$ contains the entire
two-variable simplex of degree $999974$, and
\[
                    |A|\geq\binom{999976}{2}
                        =499975500300>12!,
                                                               \tag{O-025.7}
\]
a contradiction.

It remains to treat $2\leq b\leq24$. Then
$0\leq\ell\leq11$ and
$k=26-b=24-2\ell\geq2$.
Put $R'=r-3$. The support $\Gamma_\lambda$ contains
every squarefree degree-$26$ monomial in those last
$R'$ variables. In its decomposition $\gamma=\alpha+\eta$,
the exponent $\alpha$ is squarefree and has degree at
least $k$. By downward closure, $A$ contains a
squarefree $k$-submonomial of $\gamma$.
Thus the squarefree degree-$k$ members of $A$ meet every
$26$-subset of the suffix variables. Double-counting
the incidences gives
\[
 |A|\geq
 \frac{\binom{R'}{26}}{\binom{R'-k}{26-k}}
       =\frac{\binom{R'}k}{\binom{26}k}
       \geq\frac{r^k}{2^k\,26!}.                       \tag{O-025.8}
\]
The final inequality follows because every numerator
factor in $\binom{r-3}k$ is at least $r/2$;
the additional factor $(26-k)!$ was harmlessly omitted.
Also
$T_b\geq m r^\ell/\ell!$ and $n=m r^{12}$.
It follows that
\[
 \frac{T_b|A|}{n}
       \geq\frac{r^{12-\ell}}{2^k\ell!\,26!}>1.
                                                               \tag{O-025.9}
\]
Indeed $12-\ell\geq1$ and
$2^k\ell!\,26!\leq2^{24}11!\,26!<2^{139}<r=2^{400}$.
This contradicts (O-025.5) for every remaining base order.
All cases are exhausted, proving the asserted strict
separation from every one-sided M-018 choice.

\paragraph{Other established failures remain.}
The full support (O-022.6) is contained in the present
support, and every present row cap is at least its
O-022 counterpart. Thus O-022's exact generic
degree-$26$ complexity argument still excludes every
F-018 exact support cover, while its separate dimension
arguments exclude every entire one-cut family and the
whole point-jet target.
The coordinate-box size failure also persists.
The shape meets the BIP staircase and A-025's joint gate.
The shape-specific Chow threshold is at least
$d-1=3n-1>n$, and the uniform Chow and Grenet thresholds
fail as before. Its length is $m^2$, its tail is
$26r+1999928\leq md$, and its third row is larger than two.
It remains outside every one-sided jet exclusion by the
complete argument (O-025.5)--(O-025.9).
The new two-sided comparison therefore reduces the
actual remaining gate on a shape with at least four
permanent copies. It does not produce a multiplicity
obstruction; (O-025.2) rules one out here.

\paragraph{Geometric dictionary and scope.}
The two axis ideals are finite free integral torus-weight
modules under SGA~3 I, Proposition~4.7.3, pp.~35--36.
The constant bridge and the two attached families in
M-019 give their actual universally injective coordinate
pullback. F-016 supplies the full Schur-copy support
through SGA~6 VI, Theorem~3.3, pp.~392--393, and
Macdonald I(4.3), p.~63, Appendix~A(5.3)--(5.6),
pp.~155--156. The counting argument (O-025.8)
is a positive incidence count in the suffix monomial
basis, not an alternating virtual identity.
The exact representation barrier inherited from O-022
and its border dimension barriers retain their distinct
scopes. No general saturation or unresolved evaluation-rank
comparison is inferred.

\section{Exact positive cover counts for a two-prefix profile}

\subsection{O-026: a fixed-base optimization within the existing criteria}
Let $L\geq M\geq b\geq1$, let $r\geq4$, and put
\[
 q=r-3,\qquad u=L-b,\qquad v=M-b,\qquad
 c=(v+1)(u+1)-\frac{v(v+1)}2.
\]
Consider the support
\[
 \Gamma=\left\{\gamma:
       |\gamma|\leq L,\quad
       \sum_{i=3}^r\gamma_i\leq M,\quad
       \sum_{i=4}^r\gamma_i\leq b\right\}.
                                                               \tag{O-026.1}
\]
Assume
\[
                             q>\max\{c,u+v+2\}.
                                                               \tag{O-026.2}
\]
Among nonempty finite monomial order ideals, the following
two exact minima hold:
\[
 \begin{split}
 \min_{\Gamma\subseteq A+D_b}|A|&=c,\\
 \min_{\Gamma\subseteq A+B+D_b}(|A|+|B|)&=u+v+2.
 \end{split}                                                  \tag{O-026.3}
\]
The unique one-sided optimum is
\[
                  A_0=\{x_2^a x_3^t:a+t\leq u,\ t\leq v\}.
                                                               \tag{O-026.4}
\]
Two axis ideals of degrees $u$ and $v$ attain the
two-sided minimum. These are exact positive counts for
the fixed-base support optimization in M-018 and M-019.
They are consequences and optimality statements within
those already proved criteria, not new existential
multiplicity criteria.

\begin{proof}[The one-sided optimum]
O-024 gives $\Gamma\subseteq A_0+D_b$, and its
subpartition count gives $|A_0|=c$.
Take any proposed order ideal $A$ with $|A|\leq c$.
Since there are $q>c$ suffix variables, some unit
$x_j$, $j\geq4$, is absent from $A$.
Downward closure implies that every member of $A$
has zero $j$-th exponent.

For every prefix monomial $\nu\in A_0$,
$x_j^b\nu$ belongs to $\Gamma$:
its total degree is at most $b+u=L$, its suffix from
coordinate three is at most $b+v=M$, and its final
suffix degree is $b$. In a decomposition of its exponent
as $\alpha+\eta$, with $\alpha\in A$ and $|\eta|\leq b$,
the $j$-th exponent forces $\eta=be_j$.
Thus $\alpha=\nu$, proving $A_0\subseteq A$.
Therefore $|A|\geq c$; equality forces $A=A_0$.
Any ideal with more than $c$ members cannot improve
the known optimum. This proves both the minimum and
its uniqueness.
\end{proof}

\begin{proof}[The two-sided optimum]
The two axis ideals
$A=\{1,x_2,\ldots,x_2^u\}$ and
$B=\{1,x_3,\ldots,x_3^v\}$ contain a sum covering
$A_0$, so they attain the upper bound $u+v+2$.
Now take any proposed pair with
$|A|+|B|\leq u+v+2$.
By (O-026.2) there is a suffix unit $x_j$ in neither
ideal, and no member of either ideal uses that variable.

Apply the cover to $x_j^b x_2^u$.
Its $D_b$ part must be precisely $x_j^b$, so its
remaining decomposition uses pure $x_2$ powers in
$A$ and $B$. If $a_2,b_2$ are their largest $x_2$
axis degrees, this forces $a_2+b_2\geq u$.
Likewise $x_j^b x_3^v$ gives $a_3+b_3\geq v$.
Every ideal contains its two full axis chains, sharing
only the constant monomial. Consequently
\[
 |A|+|B|\geq(a_2+a_3+1)+(b_2+b_3+1)
                              \geq u+v+2.
\]
Larger pairs cannot improve the already attained upper
bound, proving the second minimum.
\end{proof}

\paragraph{The exact counts in the positive example.}
For O-025, take
$L=1000006$, $M=1000000$, and $b=26$.
Then
\[
 \begin{split}
 u&=999980,\qquad v=999974,\\
 c&=499981500150,\qquad u+v+2=1999956.
 \end{split}                                                  \tag{O-026.5}
\]
The suffix count $r-3=2^{400}-3$ satisfies (O-026.2).
At the common fixed core size
$T=m\binom{r+11}{12}$, the optimal one-sided
matrix size is therefore $499981500151T$, whereas
the optimal two-sided size at this same base order is
$1999956T$.
O-025 has already proved that the latter fits inside
$n$ and the former does not, with actual permanent
multiplicity at least four.

For this particular example the one-sided count can
also be optimized over every even base order of M-018.
The optimum still occurs at $b=26$ with the unique
ideal (O-026.4). Indeed its size is less than
$2^{39}n$, since $T<n$ and
$499981500151<2^{39}$.
For $b\geq28$, (O-025.6) gives a lower ratio to $n$
of at least $2r/13!>2^{39}$.
For $b\leq24$, (O-025.9) gives a lower ratio greater
than $r/2^{139}=2^{261}>2^{39}$.
These cases exhaust the other orders. Enlarging the
ambient length cannot lower the core size or help
decompose monomials supported on the original variables.
This establishes the stated optimum for the M-018
architecture, without an optimality assertion for
arbitrary determinant maps or for all M-019 base orders.

\paragraph{Geometric dictionary.}
The counts in (O-026.3) are ranks of explicitly indexed
free coefficient modules over $\mathbb Z$, with their
torus-comodule structures from SGA~3 I,
Proposition~4.7.3, pp.~35--36.
Nonnegative character exponents make the missing-variable
argument exact: a monomial containing that character
cannot be cancelled inside a support decomposition.
The already proved morphisms M-018 and M-019 convert
the counts into universally injective coordinate
pullbacks. Their characteristic-zero Schur-copy step
uses SGA~6 VI, Theorem~3.3, pp.~392--393, and the
Cauchy/Schur formalism of Macdonald I(4.3), p.~63,
Appendix~A(5.3)--(5.6), pp.~155--156.
The optimization concerns these explicit coefficient
schemes and maps; neither a virtual lambda-ring identity
nor purity alone provides the comparison.

\section{Two transverse labels per matrix edge}

\subsection*{Finding O-027: the exact determinant row threshold when \(\lambda_3\leq2\)}
Let \(n\geq2\). For every outer degree \(d\) and every
\(\lambda\vdash nd\) with
\[
                    \ell(\lambda)\leq2n,\qquad \lambda_3\leq2,
                                                               \tag{O-027.1}
\]
the determinant orbit closure retains all ambient copies:
\[
                         K_n(\lambda)=a_\lambda(d[n]).
                                                               \tag{O-027.2}
\]
For the unpadded permanent of size \(n\), the same conditions give
\[
                         q^n_\lambda(d[n])=a_\lambda(d[n]).
                                                               \tag{O-027.3}
\]
These statements include all highest-weight multiplicity copies.
They follow from a single explicit construction which makes the complete
projection to transverse order two along a binary plane dominant for
both determinant and permanent in \(2n\) variables.
The construction improves O-020's sufficient determinant size
\(n\geq r-1\) to \(n\geq\lceil r/2\rceil\).
O-020 and its integral incidence calculation remain valid independently.

For \(n\geq3\), the row threshold \(2n\) in the universal assertion
(O-027.2) is sharp. The known LMR module has
\[
 \Lambda_n=\bigl(2n^3-4n^2+1,\ 2n^2-4n+1,\ 2^{\,2n-1}\bigr),
 \qquad d=2n(n-1),
                                                               \tag{O-027.4}
\]
so that \(\ell(\Lambda_n)=2n+1\) and \((\Lambda_n)_3=2\).
Theorem 2.3.1 and Section 3.2 of Landsberg--Manivel--Ressayre
give a nonzero copy of this module in the determinant ideal.
Consequently \(K_n(\Lambda_n)<a_{\Lambda_n}(d[n])\).
This sharpness statement concerns ambient determinant multiplicity;
it does not assert a permanent comparison for all \(n\).
The finite permanent endpoint certificates in A-031 supply that
additional comparison for their stated sizes.

\paragraph{The binary pencil and its quadratic coefficient space.}
Put \(N=n-1\), and use the variables
\[
                 x,y,\ z_{i1},z_{i2}\quad(1\leq i\leq N).
                                                               \tag{O-027.5}
\]
Let \(V=\bigoplus_{i=1}^{N}V_i\), where
\(V_i=\langle z_{i1},z_{i2}\rangle\), and choose distinct scalars
\(a_1,\ldots,a_n\). Write
\[
                \ell_i=y-a_i x,\qquad \Delta=\prod_{i=1}^n\ell_i.
\]
The desired coefficient space is
\[
 \mathcal J=
 \operatorname{Sym}^n\langle x,y\rangle
 \ \oplus\
 \operatorname{Sym}^{n-1}\langle x,y\rangle\otimes V
 \ \oplus\
 \operatorname{Sym}^{n-2}\langle x,y\rangle\otimes\operatorname{Sym}^2V.
                                                               \tag{O-027.6}
\]
Projection to \(\mathcal J\) discards only terms of total \(z\)-degree
at least three. Higher discarded coefficients are unrestricted.

For \(i<j<n\), set
\[
 A^0_{ij}=z_{i1}+2z_{j1},\qquad
 A^0_{ji}=z_{i2}+2z_{j2}.
\]
For \(i<n\), set
\[
                  A^0_{in}=z_{i1},\qquad A^0_{ni}=z_{i2},
\]
and set \(A^0_{ii}=\ell_i\).
Thus \(A^0\) is a completely specified matrix of linear forms.
The quadratic transverse part of its permanent is
\[
       \sum_{i<j}\frac{\Delta}{\ell_i\ell_j}A^0_{ij}A^0_{ji};
                                                               \tag{O-027.7}
\]
the corresponding determinant part is its negative.
Indeed, a permutation with exactly two off-diagonal entries is a
transposition. All cycles of length at least three disappear in
the projection (O-027.6).

\paragraph{The two-label incidence lemma.}
Let
\[
 H_n=\{(c_1,\ldots,c_n):\textstyle\sum_i c_i=0\}.
\]
For \(i<j<n\), put
\[
 U_{ij}=\langle z_{i1}+2z_{j1},\,z_{i2}+2z_{j2}\rangle,
 \qquad U_{in}=V_i.
\]
The linear map
\[
 \Psi:\bigoplus_{i<j}(U_{ij}\otimes V)
       \longrightarrow H_n\otimes\operatorname{Sym}^2V,
 \qquad
 (u\otimes v)_{ij}\longmapsto(e_i-e_j)\otimes uv
                                                               \tag{O-027.8}
\]
is surjective over every field of characteristic different from
two and three.

\begin{proof}
A functional on the target can be represented by symmetric bilinear
forms \(Q_1,\ldots,Q_n\) on \(V\), modulo adding the same form to every
\(Q_i\). Set \(Q_n=0\). A functional annihilates the image of \(\Psi\)
exactly when
\[
             (Q_i-Q_j)(u,v)=0
       \quad\text{for all }u\in U_{ij},\ v\in V.
                                                               \tag{O-027.9}
\]
The reference edges \(\{i,n\}\) first imply that \(Q_i\) kills \(V_i\)
in either argument.

For \(i<j<n\), apply (O-027.9) to
\(u=z_i+2z_j\), where the coordinates of \(z_i\in V_i\) and
\(z_j\in V_j\) agree under the displayed bases.
Taking the other argument in \(V_i\) shows that the
\(V_i\)-diagonal block of \(Q_j\) is zero.
Taking it in \(V_j\), and dividing by two, shows that the
\(V_j\)-diagonal block of \(Q_i\) is zero.
Together with the reference edges, this kills every diagonal
two-by-two block of every \(Q_i\).

Every block still possibly nonzero is indexed by three distinct
nonreference vertices. Fix \(i<j<k<n\), and denote by
\[
            A=(Q_i)_{jk},\qquad B=(Q_j)_{ik},\qquad C=(Q_k)_{ij}
\]
their two-by-two matrices. The edges \(\{i,j\}\), \(\{i,k\}\),
and \(\{j,k\}\), respectively, give
\[
                         B=2A,\qquad C=2A^{\mathsf T},
                         \qquad C=2B.
                                                               \tag{O-027.10}
\]
Hence \(A^{\mathsf T}=2A\), and transposition gives \(A=4A\).
Characteristic different from three forces \(A=0\), and then
\(B=C=0\). This treats every remaining block.
There is no triple to treat when \(N<3\); the preceding diagonal
argument already gives the same conclusion.
The dual kernel is zero, proving surjectivity.
\end{proof}

\paragraph{From binary residues to the full coefficient differential.}
The map
\[
 \operatorname{Sym}^{n-2}\langle x,y\rangle\longrightarrow H_n,
 \qquad
 f\longmapsto
 \left(\frac{f(1,a_i)}{\prod_{j\ne i}(a_i-a_j)}\right)_{i=1}^n
                                                               \tag{O-027.11}
\]
is an isomorphism. Its image is in \(H_n\) by the coefficient
of \(y^{-1}\) in the partial-fraction expansion of
\(f(1,y)/\Delta(1,y)\). The map is injective because a polynomial
of degree at most \(n-2\) vanishing at all \(n\) distinct \(a_i\)
is zero; source and target both have dimension \(n-1\).
In particular,
\[
 \frac{\Delta}{\ell_i\ell_j}\longmapsto
                    \frac{e_i-e_j}{a_i-a_j}.
                                                               \tag{O-027.12}
\]

Vary a single off-diagonal entry \(A_{ij}\) by an arbitrary
linear form \(H\in V\). The transverse-order-two part of the
derivative is exactly
\[
              \frac{\Delta}{\ell_i\ell_j}\,A^0_{ji}H
                                                               \tag{O-027.13}
\]
for the permanent and its negative for the determinant.
It has no transverse degree zero or one. Varying both
orientations of every edge, using (O-027.12), gives precisely
the image in (O-027.8), up to nonzero scalar factors.
The incidence lemma therefore supplies the whole third
summand of (O-027.6).

Next vary \(A_{ii}\) by \(H\in V\). Up to transverse order two,
the derivative is
\[
                         \frac{\Delta}{\ell_i}H.
                                                               \tag{O-027.14}
\]
There is no transverse-degree-one term in the remaining
principal minor, since a nonidentity permutation cannot have
exactly one off-diagonal entry. The \(n\) binary forms
\(\Delta/\ell_i\) form a basis in degree \(n-1\), as evaluation
at \(a_i\) shows. These variations supply the entire middle
summand, without affecting the other two.

Finally vary each diagonal entry by \(x\) or \(y\).
The transverse-degree-zero images span
\(\operatorname{Sym}^n\langle x,y\rangle\), because they are
\(x\Delta/\ell_i\) and \(y\Delta/\ell_i\).
Their possible transverse-degree-two components can be
subtracted using (O-027.13); their transverse-degree-one
components are zero. The complete differential onto
\(\mathcal J\) is consequently surjective, for both determinant
and permanent.

\paragraph{Dominance and all-copy consequence.}
The coefficients of a matrix of linear forms are affine
parameters, and projection of its determinant or permanent
to (O-027.6) is a polynomial map.
A surjective differential at \(A^0\) gives an injective
coordinate pullback: restrict source directions to a linear
right inverse of the differential, translate source and
target to their base values, and use the lowest nonzero
homogeneous part of a polynomial under a map
\(t\mapsto t+O(t^2)\).
Thus the coefficient projection is dominant.
Setting unnecessary variables to zero gives the same
conclusion for every \(r\leq2n\): coefficients using only
the retained variables are independent of the discarded
matrix-entry coefficients, so variable restriction commutes
with projection to the retained target coordinates.

For any \(\lambda\) satisfying (O-027.1), every highest-weight
vector in the degree-\(d\) ambient coordinate ring uses only
the first \(\ell(\lambda)\) variable directions.
The flag-minor support theorem F-016 says that every coefficient
generator appearing in such a vector has total exponent at
most \(\lambda_3\leq2\) in coordinates three and later.
It is therefore a polynomial in the coordinates of
\(\mathcal J\).
If the vector lies in the determinant ideal, it vanishes
on every matrix-pencil substitution into the determinant,
since such substitutions belong to the orbit closure.
The injective pullback just proved forces it to be zero.
The same reasoning applies to the unpadded permanent.
Every vector in the full highest-weight ideal multiplicity
space is covered by this argument. Consequently that space
is zero, proving (O-027.2)--(O-027.3).

\paragraph{Internal padding.}
The argument above varies all matrix-entry coefficients.
In particular, it does not claim that \(A_{11}=x\) can be
fixed in the same construction.
For the internally padded permanent, the verified O-021
padding argument instead gives the following consequence:
if \(n\geq m\), \(\ell(\lambda)\leq r\), \(\lambda_3\leq2\),
and
\[
                  m\geq
                \max\{\lceil r/2\rceil+1,\ \lambda_2+1\},
                                                               \tag{O-027.15}
\]
then
\[
                           q^m_\lambda(d[n])
                         =K_n(\lambda)=a_\lambda(d[n]).
                                                               \tag{O-027.16}
\]
Indeed, put \(s=m-1\). The new construction supplies the
full permanent coefficient projection in size \(s\), because
\(r\leq2s\). The other inequality gives \(s\geq\lambda_2\),
so F-014 stabilizes the complete highest-weight space at
base degree \(s\). The diagonal, invertible differential
transpose in O-012 preserves its transverse-order-two
coefficient support. Finally the block matrix
\(\operatorname{diag}(x,A_s)\) has fixed internal entry \(x\)
and permanent \(x\operatorname{per}(A_s)\), exactly as in
O-021. This proves all-copy equality after padding.
For one-row partitions, the nonzero pure-power coordinate
gives the same equality directly.

\paragraph{Sharp boundary for the determinant coefficient projection.}
For \(r\) variables, the target in transverse order two has dimension
\[
 D_n(r)=(n+1)+n(r-2)+(n-1)\binom{r-1}{2}.
                                                               \tag{O-027.17}
\]
The closure of the \(r\)-variable determinant-pencil image
has dimension at most \((r-2)n^2+2\), by O-010's
left-right gauge argument. At \(r=2n+1\), subtraction gives
\[
             D_n(2n+1)-\bigl((2n-1)n^2+2\bigr)=n-1>0.
                                                               \tag{O-027.18}
\]
Hence the full coefficient projection cannot be dominant
there. For larger \(r\), the difference increases.
Together with the construction, this makes \(r=2n\) the
exact maximal number of variables for this determinant
coefficient projection, within the ambient variable range.
This is a geometric statement about the whole projection;
the separate representation-theoretic sharpness in
(O-027.4) uses the known LMR ideal copy.

\paragraph{The selected integral incidence cokernel.}
The incidence map (O-027.8) is defined over \(\mathbb Z\).
With \(N=n-1\), its finite cokernel is
\[
 \operatorname{coker}\Psi_{\mathbb Z}
 \simeq
 (\mathbb Z/2)^{\,3\binom N2+3\binom N3}
       \ \oplus\
 (\mathbb Z/6)^{\,\binom N3}.
                                                               \tag{O-027.19}
\]
Here binomial coefficients are zero outside their usual range.
To see this, use the dual matrix solely to compute Smith
invariants, which are unchanged by transposition.
The reference-edge equations eliminate, by unit pivots,
every block of \(Q_i\) involving \(V_i\).
For each pair \(i<j\), the \(V_i\)-diagonal block of \(Q_j\)
is eliminated by unit pivots, while the three independent
entries of the \(V_j\)-diagonal block of \(Q_i\) have the
relations \(2q=0\).
For each triple \(i<j<k\), use the unit relations
\(B=2A\) and \(C=2B\) in (O-027.10).
The remaining four-entry block is
\[
                          2(A^{\mathsf T}-2A)=0.
\]
Its two diagonal entries contribute \(2,2\); on its two
off-diagonal entries the matrix is
\[
                    \begin{pmatrix}-4&2\\2&-4\end{pmatrix},
\]
whose Smith invariants are \(2,6\).
All pairs and triples involve disjoint remaining variables,
giving (O-027.19).
This is the cokernel of the selected incidence map;
it is not a claim that every full determinant or permanent
jet map must fail in characteristic two or three.

\paragraph{Grothendieck dictionary and base-change scope.}
The matrix pencil, determinant, permanent, coefficient
projection and incidence homomorphism are polynomial or
linear constructions over \(\mathbb Z\).
Equation (O-027.19) makes the integral limitation explicit.
After inverting six and the differences \(a_i-a_j\),
the incidence and binary interpolation maps are split
surjective, and the full differential has a linear
right inverse. The resulting lowest-term coordinate
injection survives arbitrary commutative base change
over this localized ring.
The passage to irreducible highest-weight multiplicities
is made over \(\mathbb C\), using F-016 and
characteristic-zero semisimplicity.
In Grothendieck's language, the graded coordinate algebras
and their ideals are comodules as in SGA~3, I,
Propositions~4.7.2--4.7.3, pp.~34--36.
Their Schur operations belong to the special
\(\lambda\)-structures of SGA~6, VI, Theorem~3.3,
pp.~392--393. The conclusion is certified by the actual
coordinate injection and its explicit integral localization,
not by positivity of a virtual \(K\)-class.

\section{Exact determinant multiplicities along the stretched Hessian family}

\subsection*{Finding O-028: invariant rectangle shifts of two-row cores}
Fix \(n\geq2\), put \(R=n^2\), and let
\[
 \Phi\in\mathbb C[\overline{\operatorname{GL}_R\cdot\det_n}]
\]
be a nonzero homogeneous \(\operatorname{SL}_R\)-invariant of
positive degree \(e\).
Write its polynomial highest weight as
\(\delta\mathbf1_R\), where \(\delta\in\mathbb Z_{>0}\).
Scalar weights give
\[
                              e=n\delta.
                                                               \tag{O-028.1}
\]
Such an invariant exists by the known polystability and
fundamental-invariant theorem used in M-020.
No numerical value of its degree is required below.

Let \(\sigma\vdash na\) have at most two rows, allowing the
empty partition when \(a=0\). For every integer \(j\geq0\), put
\[
               k=j\delta,\qquad
               \tau=\sigma+k\mathbf1_R,\qquad
               d=a+nk.
                                                               \tag{O-028.2}
\]
Then all the following actual and orbit multiplicities agree:
\[
 \boxed{\quad
 K_n(\tau)
   =\operatorname{sk}\bigl(\tau,(d^n)\bigr)
   =g\bigl(\tau,(d^n),(d^n)\bigr)
   =g\bigl(\sigma,(a^n),(a^n)\bigr)
   =K_n(\sigma).
 \quad}                                                       \tag{O-028.3}
\]
The two-row value on the right has the positive partition
difference formula in O-006 and F-002.
The assertion is about an actual rectangular semi-invariant
present in the determinant coordinate ring. An arbitrary
rectangle shift is not substituted for that hypothesis.

\begin{proof}[Preserving every two-row multiplicity copy]
Let \(A=\mathbb C[\overline{\operatorname{GL}_R\cdot\det_n}]\).
The determinant orbit closure is irreducible, so \(A\) is a
domain. Multiplication by \(\Phi^j\) therefore injects the
entire highest-weight space of weight \(\sigma\) into that
of weight \(\tau\), with degree changing from \(a\) to
\[
                         a+je=a+nk=d.
\]
Consequently \(K_n(\tau)\geq K_n(\sigma)\).

For the reverse bound, write \(W=U\otimes V\), where
\(\dim U=\dim V=n\). The usual determinant twist identity gives
\[
 \mathbb S_{\sigma+k\mathbf1_R}(U\otimes V)
   =\det(U\otimes V)^k\otimes\mathbb S_\sigma(U\otimes V).
\]
Since
\[
 \det(U\otimes V)^k=(\det U)^{nk}(\det V)^{nk},
\]
the multiplicity of
\((\det U)^d\otimes(\det V)^d\) on the left is exactly the
multiplicity of
\((\det U)^a\otimes(\det V)^a\) in
\(\mathbb S_\sigma(U\otimes V)\). Thus
\[
                 g(\tau,(d^n),(d^n))
                      =g(\sigma,(a^n),(a^n)).
                                                               \tag{O-028.4}
\]
This is the exact rectangular Kronecker shift identity;
it is an ordinary coefficient identity and uses no
assumed equality of orbit and closure multiplicities.

O-006 proves \(K_n(\sigma)=g(\sigma,(a^n),(a^n))\) for
every two-row \(\sigma\). The established comparison
\(K_n\leq\operatorname{sk}\leq g\), together with the
injection and (O-028.4), now bounds every term in
(O-028.3) between the same number. This proves equality.
The empty partition and \(j=0\) cases are included:
the degree-zero highest-weight space consists of constants.
\end{proof}

\paragraph{Application to the Hessian-bracket obstruction.}
For \(n\geq3\), use the A-030 notation
\[
 R=n^2,\qquad D=R+1,\qquad
 \lambda_n=(n+R(n-2)+1,\ 3,\ 2^{R-2}).
                                                               \tag{O-028.5}
\]
Set
\[
       t=(n-1)^2=D-2n,\qquad E=nt,\qquad
       h_n(s)=[z^s]\prod_{q=2}^{n}(1-z^q)^{-1}.
                                                               \tag{O-028.6}
\]
For every integer \(s\geq1\) satisfying \(\delta\mid2s\),
\[
 \boxed{\quad
       K_n(s\lambda_n)
       =\operatorname{sk}\bigl(s\lambda_n,((sD)^n)\bigr)
       =g\bigl(s\lambda_n,((sD)^n),((sD)^n)\bigr)
       =h_n(s).
 \quad}                                                       \tag{O-028.7}
\]
In particular, put
\[
                            p=\frac{\delta}{\gcd(\delta,2)}.
\]
The actual determinant multiplicity is known exactly for
every positive \(s\) divisible by \(p\). It is strictly
positive on this progression whenever \(s\geq2\).
The integer \(p\) is determined by the chosen actual
invariant; a numerical value is not asserted for general \(n\).

\begin{proof}
Padding all partitions with zeros to length \(R\), direct
subtraction gives
\[
        \lambda_n-2\mathbf1_R=(E-1,1),\qquad
        s\lambda_n=(s(E-1),s)+2s\mathbf1_R.
                                                               \tag{O-028.8}
\]
The core \(\sigma_s=(s(E-1),s)\) has size \(nst\), so
its outer degree is \(st\). The exact two-row formula
O-002, with rectangle height \(n\) and width \(st\), yields
\[
                    K_n(\sigma_s)
                   =g(\sigma_s,((st)^n),((st)^n))
                   =h_n(s).
                                                               \tag{O-028.9}
\]
The width bound is inactive because \(st\geq s\).
If \(\delta\mid2s\), multiply the complete core space by
\(\Phi^{2s/\delta}\). Its degree increases by
\[
                     e\,\frac{2s}{\delta}=2sn,
\]
and the total degree is \(st+2sn=sD\), exactly the degree
required for \(s\lambda_n\). Apply (O-028.3) to obtain
(O-028.7).

Finally \(h_n(1)=0\). For \(n\geq3\), both parts two and
three occur in (O-028.6), so every integer \(s\geq2\)
has a representation by those parts and \(h_n(s)>0\).
This proves the positivity claim on the stated progression.
\end{proof}

\paragraph{Orbit multiplicity at every dilation.}
The ordinary and symmetric orbit coefficients in fact satisfy
\[
 \operatorname{sk}\bigl(s\lambda_n,((sD)^n)\bigr)
   =g\bigl(s\lambda_n,((sD)^n),((sD)^n)\bigr)
   =h_n(s)\qquad(s\geq1),
                                                               \tag{O-028.9a}
\]
without the divisibility condition \(\delta\mid2s\).
For the ordinary coefficient this is the determinant twist
in (O-028.8)--(O-028.9).
For the symmetric coefficient, transpose on
\(\mathbb C^n\otimes\mathbb C^n\) has determinant
\((-1)^{n(n-1)/2}\).
The determinant-line twist of exponent \(2s\) therefore
has transpose scalar one.
Consequently the same identification of rectangular
multiplicity spaces preserves the transpose involution
and its symmetric subspace.
Its two-row core has symmetric coefficient equal to the
ordinary coefficient by O-006, proving (O-028.9a).
Thus the orbit multiplicity is already positive at every
\(s\geq2\); an actual invariant section is still needed
for the closure equalities proved above and below.

\paragraph{Every sufficiently large dilation, and a genuine support hole.}
An additional known theorem on the invariant degree monoid
strengthens the progression conclusion to all sufficiently
large integers. For every fixed \(n\geq3\), there is an integer
\(s_0(n)\) such that
\[
                       K_n(s\lambda_n)=h_n(s)
                  \qquad\text{for every }s\geq s_0(n).
                                                               \tag{O-028.11}
\]
No numerical bound on \(s_0(n)\) is asserted.
The exact ordinary and symmetric coefficients in
(O-028.7) equal the same value for all these \(s\).
In particular the actual determinant support monoid
\[
                        S_n=\{\mu:K_n(\mu)>0\}
\]
has the explicit relative-saturation hole
\[
                           \lambda_n\in
                         \operatorname{Sat}(S_n)\setminus S_n.
                                                               \tag{O-028.12}
\]
Here relative saturation uses the group generated by \(S_n\);
membership in that group is proved below, rather than inferred
from the existence of a single positive stretch.

\begin{proof}
Let \(\mathcal E_n\) be the monoid of degrees of nonzero
\(\operatorname{SL}_{n^2}\)-invariants in the determinant
coordinate ring.
The additional primary inputs are
B\"urgisser--Ikenmeyer, \emph{Fundamental invariants of orbit
closures}, arXiv:1511.02927v2:
\begin{itemize}
\item Theorem 2.5, printed p.~5, gives the determinant
stabilizer period
\[
 \alpha_n=
 \begin{cases}
  1,&n\bmod4\in\{0,1\},\\
  2,&n\bmod4\in\{2,3\};
 \end{cases}
\]
\item Definition 3.1 and Theorem 3.4, printed p.~8, say
that the invariant degree monoid generates the group
\[
                          \operatorname{gp}(\mathcal E_n)
                                    =n\alpha_n\mathbb Z.
\]
\end{itemize}
The theorem's polystability hypothesis is supplied by
Corollary 2.9, pp.~6--7, as already used in M-020.
Thus
\[
                    M_n=\frac1{n\alpha_n}\mathcal E_n
\]
is an additive submonoid of \(\mathbb Z_{\geq0}\) whose
generated group is \(\mathbb Z\).

For completeness, this condition implies that \(M_n\)
contains all sufficiently large integers.
Choose any positive \(h\in M_n\).
The residues of \(M_n\) in the finite cyclic group
\(\mathbb Z/h\mathbb Z\) form an additive submonoid,
and hence a subgroup: every element has its inverse
among its nonnegative multiples.
Since \(M_n\) generates \(\mathbb Z\), this subgroup is
all of \(\mathbb Z/h\mathbb Z\).
Choose \(b_j\in M_n\) for each residue \(j\).
Every integer \(u\geq\max_j b_j\), in its corresponding
residue class, is \(b_j+vh\) for some \(v\geq0\) and
therefore belongs to \(M_n\).
Let \(c_n\) be any bound with this property.

Because \(\alpha_n\) is one or two, \(2s/\alpha_n\)
is an integer for every \(s\).
If
\[
                  s\geq
             \max\{2,\lceil\alpha_n c_n/2\rceil\},
\]
then \(2s/\alpha_n\in M_n\).
Hence the determinant coordinate ring contains a
nonzero invariant \(\Psi_s\) of degree \(2sn\), whose
highest weight is \(2s\mathbf1_{n^2}\).
Apply the invariant-shift lemma (O-028.3) to
\(\Psi_s\) and the same two-row core
\(\sigma_s=(s(E-1),s)\).
It gives (O-028.11), with strictly positive value because
\(s\geq2\).
There is no need to express \(\Psi_s\) as a power of one
fixed fundamental invariant.

Choose \(s\) large enough that both \(s\lambda_n\) and
\((s+1)\lambda_n\) have positive determinant multiplicity.
Their difference shows
\[
                   \lambda_n=(s+1)\lambda_n-s\lambda_n
                                  \in S_n-S_n.
\]
Some positive multiple of \(\lambda_n\) is in \(S_n\),
so this proved group membership implies
\(\lambda_n\in\operatorname{Sat}(S_n)\).
Finally A-030 gives \(K_n(\lambda_n)=0\).
This proves the hole assertion (O-028.12).
\end{proof}

A-030 also proves that the actual unpadded permanent
multiplicity at \(s\lambda_n\) is positive for every \(s\geq1\),
by powers of its nonzero Hessian-bracket highest-weight section.
Thus the occurrence obstruction present at \(s=1\) disappears
for every sufficiently large dilation, with the exact
determinant multiplicity given by (O-028.11).
The argument uses the special two-row core and the
determinant's invariant period; it is not a general
Kronecker or determinant-support saturation theorem.

The current status is stronger than this invariant-based
argument alone. A-034 proves, for every \(n\geq3\) and
every integer \(s\geq1\),
\[
 q^n_{s\lambda_n}(s(n^2+1)[n])
       \geq p_n(s)>h_n(s)=K_n(s\lambda_n),
 \qquad
 p_n(s)=[z^s]\prod_{i=1}^{n}(1-z^i)^{-1}.
\]
The determinant equality follows by multiplying the
entire binary highest-weight space by a nonzero
Hessian section; it requires no invariant-period
threshold. Hence both actual multiplicities are
positive at every \(s\geq2\), and the strict gap
persists at all dilations. O-030 extends that
construction to a three-parameter flat-tail family.
For \(n=3\), F-028's injective binary
permanent--Hessian coefficient map supplies a stronger
permanent lower bound. The invariant-shift lemma and
numerical-monoid proof above remain valid independent
tools; their ray consequences are superseded by these
all-integer formulas.

\paragraph{A concrete even-size specialization.}
The known invariant computation for \(n=4\) makes one
progression numerical.
B\"urgisser--Ikenmeyer, \emph{Fundamental invariants of orbit
closures}, arXiv:1511.02927v2, Theorem 3.15,
printed p.~13, gives the general minimal-degree bound.
The paragraph following Proposition 3.28, printed p.~18,
reports the computation
\[
                            P_{4,16}(\det_4)\ne0.
\]
Together these statements give a nonzero invariant of
degree \(e=16\), and hence \(\delta=4\).
This nonvanishing is a used published computation, not
recomputed or relabeled as new here.
The original PDF is
\url{https://arxiv.org/pdf/1511.02927v2}, retained locally
as \path{external_sources/BI_fundamental_invariants.pdf}.

Here \(D=17\) and
\(\lambda_4=(37,3,2^{14})\).
For every integer \(u\geq1\), (O-028.7) becomes
\[
 \begin{split}
 K_4\bigl(2u(37,3,2^{14})\bigr)
   &=[z^{2u}]\frac1{(1-z^2)(1-z^3)(1-z^4)}\\
   &=\sum_{v=0}^{\lfloor u/3\rfloor}
                    \left(\left\lfloor\frac{u-3v}{2}\right\rfloor+1\right).
 \end{split}                                                   \tag{O-028.10}
\]
For the second equality, a solution to
\(2a+3b+4c=2u\) has \(b=2v\);
then \(a+3v+2c=u\), leaving exactly the displayed
number of choices for \(c\).
The first values for \(u=1,2,3,4,5,6\) are
\[
                              1,\ 2,\ 3,\ 4,\ 5,\ 7.
\]
These are actual determinant-closure multiplicities,
not merely ordinary rectangular coefficients.

\paragraph{Grothendieck dictionary and the integral limitation.}
The proof uses multiplication by a genuine nonzero section
of a one-dimensional representation, followed by the
actual coordinate-ring upper bound.
The determinant-twist identity is the Schur-functor
identity obtained by tensoring with a determinant line;
the special \(\lambda\)-structure is SGA~6, VI,
Theorem~3.3, pp.~392--393, with the functorial decomposition
in Macdonald, I, Appendix A (5.3)--(5.6), pp.~155--156.
Highest-weight products and rectangular characters are
the coordinate-comodule and grading constructions of
SGA~3, I, Propositions~4.7.2--4.7.3, pp.~34--36.
These tensor and determinant identities have integral
models. The actual section \(\Phi\), its nonvanishing on
the determinant, the irreducible multiplicity comparison,
and the exact assertions above are made over
\(\mathbb C\).
No chosen integral normalization of \(\Phi\) is assumed
nonzero on every special fiber.
The result uses neither normality nor an inference from
positivity in a virtual Grothendieck ring.


\section{What is, and is not, supplied by positivity}

Brion, \emph{Positivity in the Grothendieck group of complex flag
varieties}, Theorem~1, asserts that for a subvariety $Y$ with rational
singularities the coefficient $b_w$ in
$[\mathcal O_Y]=\sum_w b_w[\mathcal O_{X_w}]$ satisfies
$(-1)^{\operatorname{codim}X_w-\operatorname{codim}Y}b_w\geq0$.
Its application to products of Schubert structure sheaves is explained
immediately after that theorem (and uses Lemma~2). These are the
alternating positive structure constants represented in type A by
Grothendieck polynomials.

There are three separate operations here. Induction multiplication on
$\bigoplus_N R(S_N)$ has Littlewood--Richardson constants. Tensor
multiplication inside a fixed $R(S_N)$ has Kronecker constants.
Multiplication of coherent-sheaf classes on a flag variety has
Grothendieck-polynomial structure constants. Equality of these
products, or a map preserving their distinguished bases and their
stretching, has not been established by a $\lambda$-ring identity.
In particular a proposed general saturation implication from special
$\lambda$ identities is refuted by O-002: those identities hold in
the representation rings that exhibit the holes.  A proposed
implication from normality or positive multiplication alone is refuted
by O-003. Neither observation rules out an additional geometric
construction with substantially stronger hypotheses in a restricted
family. The missing extra statement must be supplied and proved.

The positive basis theorem actually established by the supplied GCT IV
is Theorem~15.21(vii), with the explicit rule in Theorem~18.2, for
two factors with at most two rows. The full compatible basis programme
is still stated there as Conjecture~19.1, and Section~19.2 explains
the dimension failure of the direct construction beyond that case.
It cannot be promoted to a proved general geometric positivity theorem
merely by identifying exterior powers in $K_0$.

The exact known occurrence conclusions in the GCT range are:
Ikenmeyer--Panova Theorem~1.6, $n>3m^4$, gives
$g_{n,d}(\lambda)=0\Rightarrow q^m_\lambda(d[n])=0$;
B\"urgisser--Ikenmeyer--Panova Theorem~1.4, $n\geq m^{25}$, gives
$q^m_\lambda(d[n])>0\Rightarrow K_n(\lambda)>0$.
The latter does not assert $q^m_\lambda\leq K_n$.
O-005--O-007 prove the latter inequality in the requested hook and
two-row families (and extend it to three-row outputs), but supply
no general multiplicity obstruction and no separation of complexity
classes.

\section{Geometric placement of the findings}

Finite symmetric groups are constant finite \'etale group schemes
over $\mathbb Z$ (SGA~3, Expos\'e~I, Section~4.1), and their
representations are comodules by SGA~3, I, Proposition~4.7.2.
Passing to $K_0$ over a characteristic-zero field turns the symmetric,
exterior, and Schur functors used above into the usual operations on
locally free representations. The special $\lambda$ identities live
in SGA~6, Expos\'e~V, Definition~2.4 and identities~(2.4.2), and
the realization for locally free sheaves is Expos\'e~VI,
Theorems~3.2--3.3 (pp.~391--393). The exterior filtration behind the
direct-sum and exact-sequence formulas is V, Lemma~2.2.1
(pp.~315--316). The $\gamma$-filtration is V, 3.10,
equation~(3.10.1), p.~331. None of these operations makes an arbitrary
virtual Burnside class effective; the counterexample in cote~133,
scan~019 must be corrected before making the translation.

O-001 and O-002 state multiplicities in
$K_0(\operatorname{Rep}_{\mathbb Q}S_N)$ and, through Schur--Weyl
duality, in the restriction of $S_\lambda(E\otimes F)$ along the
group-scheme homomorphism
$\mathrm{GL}(E)\times\mathrm{GL}(F)\to\mathrm{GL}(E\otimes F)$.
This is the Cauchy/Schur dictionary relevant to cote~133, scan~028.
O-003 has the more explicit integral graded-affine-scheme realization
already given above. O-004 is a statement about invariant summands
of Schur functors of the adjoint bundle on the characteristic-zero
fiber. O-005 is the exterior-power calculation in that exact tensor
category.

O-006.a is a statement about equivariant coherent quotients of the
coordinate algebra of an affine space of sections. O-006.b and
O-007 apply it to the scheme-theoretic orbit closure of an integral
form: the determinant and its substitution morphism are defined
over $\mathbb Z$, and their characteristic-zero fibers give the
orbit closures above. Semisimplicity and the displayed irreducible
decompositions are used only after base change to $\mathbb Q$ or
$\mathbb C$, not integrally and not in all residue characteristics.
O-007 also has a concrete Picard/Quot interpretation: Beauville's
matrix gives a resolution
$0\to\mathcal O_{\mathbb P^2}(-2)^n\to
\mathcal O_{\mathbb P^2}(-1)^n\to L\to0$ on a plane curve,
with $L$ a point of its Picard scheme outside the theta locus and
the final arrow a coherent quotient. The Hilbert/Picard/Quot
constructions provide the parameter spaces, while the determinantal
and multiplicity conclusions are the arguments of this run and the
explicit cited theorem, not statements attributed retrospectively
to Grothendieck. For the exact Quot representability statement, see
Grothendieck, FGA, Bourbaki expos\'e~221, \emph{Techniques de
construction en g\'eom\'etrie analytique, IV: Les sch\'emas de Hilbert},
Theorem~3.1, p.~259, and Theorem~3.2, p.~260, for a fixed Hilbert
polynomial; the identification of the Hilbert functor with the Quot
functor of $\mathcal O_X$ follows Proposition~3.8 on p.~265.
Remark~3.9, p.~266, explains the reduction to projective schemes over
$\mathbb Z$. O-008 is a statement about the homogeneous ideal of the
scheme-theoretic image of the same integral substitution morphism
after base change to $\mathbb C$, with its dimension bound proved
directly by the displayed parametrization.
O-009 is the invariant-multiplicity statement on the characteristic-zero
fiber obtained from the stated positivity and plethysm theorems.
O-010 is a dimension constraint on containment of a linear affine
subscheme in a scheme-theoretic image; its proof is an explicit
coordinate-ring transcendence-degree calculation, so requires no
unproved normality or purity property of that image.
